\documentclass[11pt]{article}
\usepackage[a4paper,margin=19mm]{geometry}
\usepackage{fontspec}
\setmainfont{lmroman10-regular.otf}
\usepackage{amsmath,amssymb,amsthm,mathtools,amscd,xcolor,graphicx,hyperref}
\hypersetup{hidelinks}
\numberwithin{equation}{section}
\newtheorem{theorem}{Theorem}[section]
\newtheorem{lemma}[theorem]{Lemma}
\newtheorem{proposition}[theorem]{Proposition}
\newtheorem{corollary}[theorem]{Corollary}
\newtheorem{claim}[theorem]{Claim}
\newtheorem{fact}[theorem]{Fact}
\newtheorem{conjecture}[theorem]{Conjecture}
\theoremstyle{definition}
\newtheorem{definition}[theorem]{Definition}
\newtheorem{example}[theorem]{Example}
\newtheorem{remark}[theorem]{Remark}
\newtheorem{question}[theorem]{Question}
\theoremstyle{plain}
\newcommand{\C}{\mathbb C}
\newcommand{\R}{\mathbb R}
\newcommand{\Q}{\mathbb Q}
\newcommand{\Z}{\mathbb Z}
\newcommand{\N}{\mathbb N}
\newcommand{\eee}{\mathrm e}
\newcommand{\cclass}[1]{\textsf{#1}}
\newcommand{\cproblem}[1]{\textsc{#1}}
\newcommand{\ie}{i.\,e.}
\newcommand{\Ie}{I.\,e.}
\newcommand{\eg}{e.\,g.}
\newcommand{\Eg}{E.\,g.}
\newcommand{\phd}{Ph.\,D.}
\newcommand{\msc}{M.\,S.}
\newcommand{\bsc}{B.\,S.}
\newcommand{\expref}[2]{\hyperref[#2]{#1~\ref{#2}}}
\newcommand{\expeqref}[2]{\hyperref[#2]{#1~\eqref{#2}}}
\newcommand{\exprefsilent}[2]{\hyperref[#2]{#1}}
\providecommand{\etal}{et~al.}
\usepackage{eprint}
\usepackage{bbm}
\usepackage[boxruled, noend, noline]{algorithm2e}

%
\newcommand{\eps}{\varepsilon}

\newcommand{\ord}[2][th]{\ensuremath{{#2}^{\mathrm{#1}}}}
\newcommand{\ceil}[1]{\left\lceil {#1} \right\rceil}
\DeclareMathOperator{\E}{\hbox{\bf E}}


\newcommand{\bb}{\ensuremath{\mathbf{b}}}
%
%
%
\newcommand{\bS}{\ensuremath{\mathbf{S}}}
\newcommand{\bT}{\ensuremath{\mathbf{T}}}
%
\newcommand{\cA}{\ensuremath{\mathcal{A}}}
\newcommand{\cC}{\ensuremath{\mathcal{C}}}
\newcommand{\cD}{\ensuremath{\mathcal{D}}}
\newcommand{\cP}{\ensuremath{\mathcal{P}}}
\newcommand{\cT}{\ensuremath{\mathcal{T}}}

\newcommand{\hypercube}{\ensuremath{\{0,1\}^\dd}}
\newcommand{\totcube}{\ensuremath{m}}
\newcommand{\dir}{\mathsf{dir}}
\newcommand{\Yes}{\ensuremath{\mathcal{D}^+}}
\newcommand{\No}{\ensuremath{\mathcal{D}^-}}
\newcommand{\Far}{\ensuremath{\mathsf{FAR}}}
\newcommand{\FarNo}{\ensuremath{\widehat{\mathcal{D}}^-}}
\newcommand{\lbconst}{\ensuremath{\delta_0}}
\newcommand{\order}{\ensuremath{<_\dec}}
\newcommand{\goodS}{\ensuremath{E_S}}
\newcommand{\coord}[2]{t^{#1}_{#2}}
\newcommand{\capt}{\mathtt{cap}}
\newcommand{\mpot}{\mathcal{A}_{\text{dir}}}
\newcommand{\dist}{\mathsf{dist}}
\newcommand{\dd}{\ensuremath{{d'}}}
\newcommand{\dec}{\mathsf{val}}
\newcommand{\bin}{\mathsf{bin}}
\newcommand{\sgn}{\mathsf{sgn}}
\newcommand{\supp}{\mathsf{supp}}

\newcommand{\Sec}[1]{\expref{Section}{sec:#1}}
\newcommand{\Alg}[1]{\expref{Algorithm}{alg:#1}}
\newcommand{\Thm}[1]{\expref{Theorem}{thm:#1}}
\newcommand{\Lem}[1]{\expref{Lemma}{lem:#1}}
\newcommand{\Clm}[1]{\expref{Claim}{clm:#1}}
\newcommand{\Def}[1]{\expref{Definition}{def:#1}}
\newcommand{\Step}[1]{\expref{Step}{step:#1}}


%
%
%
\title{Omega(d log d) nonadaptive queries for real-valued unateness testing}
\author{Roksana Baleshzar, Deeparnab Chakrabarty, Ramesh Krishnan S. Pallavoor, Sofya Raskhodnikova, C. Seshadhri}
\date{doi:10.4086/toc.2020.v016a003; CC BY 3.0}
\begin{document}
\maketitle
\noindent\textit{Editorial note.} Complete original mathematical text, author algorithms and proofs, and original published bibliography retained. Native frontmatter/biographies remain in raw source/publication PDF assets. Existing article interface renders ToC notation for XeLaTeX only.\par\medskip
\noindent\textit{Editorial note.} One general primary object only. Introductory restatements, source-local special cases and direct application consequences remain complete context. Exact original CC BY3.0 and model/field/parameter conditions retained.\par\medskip
\noindent\textit{Editorial note.} Authors own mathematics; no proof review, mathematical repair, added assumptions, content hashes or cross-version comparison.\par\medskip
\noindent\textit{Allowed standard prerequisites (editorial).}\begin{itemize}
\item Original model conventions and cited standard tools. Original sharing, matrix-order or testing model, parameter restrictions and probability/error conditions retained literally.. Complete source-local core argument retained without collector reproof or independent certification.
\end{itemize}\medskip
\section{Introduction}
%
We study the problem of testing whether a given real-valued function $f$ on domain $[n]^d$, where $n,d\in\N,$ is unate.
A function
$f:[n]^d \to \R$
is \emph{unate} if for every coordinate $i\in [d]$, the function is either nonincreasing in the $\ord{i}$ coordinate or nondecreasing in the $\ord{i}$ coordinate.
Unate functions naturally generalize monotone functions, which are nondecreasing in all coordinates, and $\bb$-monotone functions, which have a particular direction in each coordinate (either nonincreasing or nondecreasing), specified by a bit-vector $\bb\in \{0,1\}^d$. More precisely,
%
 a function is $\bb$-monotone if it is nondecreasing in coordinates $i$ with $\bb_i = 0$ and nonincreasing in the other coordinates. Observe that a function $f$ is unate iff there exists some $\bb\in \{0,1\}^d$ for which $f$ is $\bb$-monotone.

A \emph{tester}~\cite{RS96,GGR98} for a property $\cal P$ of a function $f$ is an algorithm that gets a distance parameter $\eps\in(0,1)$ and query access to $f$. The (relative) distance from a function $f$ to a property $\cal P$ is the smallest fraction of values of $f$ that must be modified to make $f$ satisfy $\cal P$. A function $f$ is $\eps$-far from $\cal P$ if the distance from $f$ to $\cal P$ is at least $\eps$. A tester for $\cal P$ has to accept with probability at least $2/3$ if $f$ has property $\cal P$ and reject with probability at least $2/3$ if $f$ is $\eps$-far from $\cal P$. A tester has \emph{one-sided error} if it always accepts a function satisfying $\cal P$; it has \emph{ two-sided
error} otherwise.
A \emph{nonadaptive} tester makes all its queries at once, whereas an \emph{adaptive} tester can make queries after seeing answers to previous queries.

Testing various properties of functions, including
monotonicity
(see, \eg, \cite{GGLRS00, Ras99, DGLRRS99, EKKRV00, LR01, FLNRRS02, Ras03, Fis04, HK07, AC06, HK08, BRW05, BBM12, BGJRW12, BCGM12, CS13, BlaRY14, BerRY14, CS14, CS16, CDJS17, CST14, CDST15, KMS18, BB15, BB16, DRTV16, PRV17, CWX17, B18, BCS18, CS19, BCS20, PRW20} and recent surveys~\cite{Ras10, Ras16, C16a}),
the Lipschitz property~\cite{JR13, CS13, DJRT13, BlaRY14, AJMR16}, bounded-derivative properties~\cite{CDJS17}, linearity~\cite{BLR93, BCHKS96, BCLR08, KLX10, RR16}, submodularity~\cite{PRR03, RY13, SV14, BB17}, and unateness~\cite{GGLRS00,KS16}, has been studied extensively over the past two decades. Even though unateness testing was initially discussed in the seminal paper
by Goldreich et al.~\cite{GGLRS00},
which was one of the first papers that study property testing, 
%
relatively little was known about testing this property. All previous work on unateness testing focused on the special case of Boolean functions on domain $\{0,1\}^d$. The domain $\{0,1\}^d$ is called the \emph{hypercube} and the more general domain $[n]^d$ is called the \emph{hypergrid}.
Goldreich et al.~\cite{GGLRS00}
provided 
%
%
%
an
%
$O({d^{3/2}}/{\eps})$-query nonadaptive tester for unateness of Boolean functions on the hypercube.
More than a decade later, Khot and Shinkar~\cite{KS16} improved the query complexity
to $O((d\log d)/\eps)$, albeit with an adaptive tester.

In this paper, we improve upon both these works, and our results hold for a more general class of functions.
Specifically, we show that unateness of real-valued functions on hypercubes can be tested nonadaptively with $O(( d /\eps) \log (d/ \eps))$ queries and adaptively with $O(d/\eps)$ queries. More generally, we describe 
%
%
an
%
$O((d/\eps) \cdot(\log (d/\eps) + \log n))$-query nonadaptive tester and 
%
%
an
%
$O(({d\log n})/\eps)$-query adaptive tester for unateness of real-valued functions over hypergrids.

In contrast to the state of knowledge for unateness testing, the complexity of testing monotonicity of real-valued functions over the hypercube and the hypergrid has been resolved. For constant distance parameter $\eps$, it is known to be $\Theta(d\log n)$. Moreover, this bound holds for all \emph{bounded-derivative} properties~\cite{CDJS17}, a large class that includes $\bb$-monotonicity and some properties quite different from monotonicity, such as the Lipschitz property. Amazingly, the upper bound for all these properties is achieved by the same simple and, in particular, nonadaptive tester.
Even though proving lower bounds for adaptive testers has been challenging in general, a line of work, starting from Fischer~\cite{Fis04} and including \cite{BBM12,CS14,CDJS17}, has established that adaptivity does not help for this large class of properties. Since unateness is so closely related, it is natural to ask whether the same is true for testing unateness.

We answer this in the negative: we prove that any nonadaptive unateness tester of real-valued functions over the hypercube (for some constant distance parameter) must make $\Omega(d\log d)$ queries.
More generally, it needs $\Omega(d(\log d+\log n))$ queries for the hypergrid domain. These lower bounds complement our algorithms, completing the picture for testing unateness of real-valued functions.
{}From a property testing standpoint, our results establish that unateness is different from
monotonicity and, more generally, any derivative-bounded
property.

\subsection{Formal statements and technical overview}
Our testers are summarized in the following theorem, stated for functions over the hypergrid domains. (Recall that the hypercube is a special case of the hypergrid with $n=2$.)

\begin{theorem}\label{thm:main-hg}
	Consider functions $f:[n]^d \to \R$ and a distance parameter $\eps\in (0,1/2)$.
	\begin{enumerate}
        \item\label{item:nonadaptive} There is a nonadaptive unateness tester that makes
          \[
            O\left(\frac d \eps\left(\log\frac d\eps + \log n\right)\right)
          \]
          queries.\footnote{For many properties, when the domain is
            extended from the hypercube to the hypergrid, testers
            incur an extra multiplicative factor of $\log n$ in the
            query complexity. This is the case for our adaptive
            tester.  However, the complexity of nonadaptive unateness
            testing (for constant $\eps$) is
            $\Theta(d(\log d + \log n))$ rather than
            $\Theta(d\log d\log n)$.}
		
		\item\label{item:adaptive} There is an adaptive unateness tester  that makes $O(({d\log n})/\eps)$ queries.
	\end{enumerate}
	Both testers have one-sided error.
\end{theorem}

Our main technical contribution is the proof that the extra $\Omega(\log d)$ is needed for nonadaptive testers.
This result demonstrates a gap between adaptive and nonadaptive unateness testing.

\begin{theorem}\label{thm:non-adap-lb-1}
Any nonadaptive unateness tester (even with two-sided error) for real-valued functions $f:\{0,1\}^d \to \R$ with
a %
distance parameter $\eps = 1/8$  must make $\Omega(d\log d)$ queries.
\end{theorem}


The lower bound for adaptive testers is an easy adaptation of the monotonicity lower bound in~\cite{CS14}.
We state this theorem for completeness and prove it in \Sec{adap-lb}.
\begin{theorem}\label{thm:adap-lb}
  Any unateness tester  for functions $f:[n]^d \to \R$ with 
  a %
  distance parameter $\eps \in (0,1/4)$ must make
  \[
    \Omega\left(\frac{d \log n}{\eps}-\frac{\log 1/\eps}{\eps}\right)
  \]
  queries.
\end{theorem}

Theorems~\ref{thm:non-adap-lb-1} and~\ref{thm:adap-lb} directly imply that our nonadaptive tester is optimal for constant $\eps$, even for the hypergrid domain.
The following theorem is proved in
\Sec{na-lb-hg}.

\begin{theorem}\label{thm:na-lb-hg}
Any nonadaptive unateness tester (even with two-sided error) for real-valued functions $f:[n]^d \to \R$ must make $\Omega(d(\log n + \log d))$ queries.
\end{theorem}


\subsubsection{Overview of techniques}\label{sec:intro-tech}
We first consider the hypercube domain. For each $i\in[d],$ an \emph{$i$-edge} of the hypercube is a pair $(x,y)$ of points in $\{0,1\}^d$, where $x_i=0,y_i=1$, and $x_j=y_j$ for all $j\in([d] \setminus \{i\})$. Given an input function $f:\{0,1\}^d\to\R$, we say an $i$-edge $(x,y)$ is \emph{increasing} if $f(x)<f(y)$, \emph{decreasing} if $f(x)>f(y),$ and \emph{constant} if $f(x)=f(y)$.

Our nonadaptive unateness tester on the hypercube uses the work investment strategy from~\cite{BerRY14} (also refer to Section 8.2.4 of Goldreich's book~\cite{Go-book}) to ``guess'' a good dimension where to look for violations of unateness (specifically, both increasing and decreasing edges).  For all $i\in[d]$, let $\alpha_i$ be the fraction of the $i$-edges that are decreasing, $\beta_i$ be the fraction of the $i$-edges that are increasing, and $\mu_i = \min(\alpha_i,\beta_i)$. The dimension reduction theorem from~\cite{CDJS17} implies that if the input function is $\eps$-far from unate, then the average of $\mu_i$ over all dimensions is at least $\eps/({4d})$. If the tester knew which dimension had $\mu_i=\Omega(\eps/d)$, it could detect a violation with high probability by querying the endpoints of $O(1/\mu_i)=O(d/\eps)$ uniformly random $i$-edges.
Performing this test for every dimension would result in query complexity $\Theta(d^2/\eps)$. The work investment strategy allows us to achieve query complexity $O((d/\eps) \log (d/\eps))$ by repeatedly choosing a uniformly random dimension and investing a specific number of queries in trying to find a violation of unateness in that dimension. It proceeds in $\Theta(\log (d/\eps))$ stages, doubling, in every stage, the number of queries invested in each dimension. Intuitively, when all violations are in one dimension, a tester has to try $\Theta(d)$ dimensions before it finds the bad one, but needs only $\Theta(1/\eps)$ queries in the bad dimension. In contrast, when all dimensions have $\mu_i=\eps/(4d)$, only a constant number of attempts are needed to find a dimension with violations, but $\Theta(d/\eps)$ queries in one dimension are required to detect a violation. (In this case, sampling $\Theta((d/\eps) \log (d/\eps))$ uniformly random edges would not be enough to detect a violation.) The work investment strategy allows us to interpolate between these scenarios.

%
%
%
%

With adaptivity, this search  through $\Theta(\log (d/\eps))$ different scenarios is not required.
A pair of queries in each dimension detects dimensions with many non-constant edges, and the algorithm focuses
on finding violations in those dimensions. This leads to the query complexity of $O(d/\eps)$, removing the $\log (d/\eps)$ factor.

It is relatively easy to extend (both adaptive and nonadaptive) testers from hypercubes to hypergrids by incurring an extra factor of $\log n$ in the query complexity. The role of $i$-edges is now played by \emph{$i$-lines}. An $i$-line is a set of $n$ domain points that differ only in coordinate $i$. The domain $[n]$ is called a line.
Monotonicity on the line (a.k.a.\ sortedness) was one of the first properties studied in the context of property testing~\cite{EKKRV00}. What we need is a nonadaptive tester for sortedness that has 1-sided error and is, in addition, \emph{proximity oblivious}: that is, it rejects a function $f:[n]\to\R$ with probability proportional to the distance from $f$ to monotonicity. Proximity-oblivious testers (POTs) were defined by Goldreich and Ron~\cite{GR11}. There are several POTs for sortedness that make $O(\log n)$ queries: the \emph{tree tester}~\cite{EKKRV00}, the spanners-based tester~\cite{BGJRW12, Ras10}, and the \emph{power of 2} tester~\cite{CS13}.
Such testers can be easily modified to work for unateness on the line and to output not just an accept/reject decision, but whether they found a pair of points on which the function is strictly increasing and, similarly, whether they found a pair of points on which the function is strictly decreasing (see \Sec{hypergrids} for details).
To generalize our unateness testers from the hypercube to the hypergrid domains, instead
of sampling a random $i$-edge, we sample a random $i$-line $\ell$ and run a modified POT for unateness on 
%
the restriction $f_{|\ell}$ of function $f$ to the line $\ell$.
This ``direct generalization'' is optimal for adaptive testers, but, interestingly, not for nonadaptive testers.

Intuitively, for nonadaptive testers, the ``direct generalization'' works with only $O(({d\log n})/\eps)$ queries when there are enough lines that are, cumulatively, sufficiently far from unate. However, it could happen that a function on the hypergrid is far from unate, but does not have any lines that are far from unate: the distance from unateness arises because some lines are far from monotone and other lines in the same dimension are far from 
%
%
\emph{antitone} %
(that is, %
far from nonincreasing).  %
%
We prove that each function $f$ on the line that is $\eps$-far from monotone, but is not $\eps/2$-far from unate, is strictly 
%
%
decreasing
on
an  %
$\eps/4$ fraction of pairs in $[n]$. Symmetrically, if a function is far from 
%
%
antitone %
but close to unate, it is strictly 
%
increasing 
%
on a large fraction of pairs. The dimension reduction allows us to use this statement to show that if the ``direct generalization'' does not work well, then, intuitively, the average dimension has many pairs on which $f$ is increasing and many pairs on which $f$ is decreasing. We again use
the %
work investment strategy to get a tester for this case that has the same complexity as our nonadaptive tester for the hypercube.
The resulting nonadaptive complexity (for constant $\eps$) is $O(d(\log d + \log n))$, which we show is optimal.

\paragraph{The nonadaptive lower bound.} Our most significant finding
is the $\log d$ gap in the query complexity between adaptive and
nonadaptive testing of unateness.  By techniques from previous
work~\cite{Fis04,CS14}, it suffices to prove lower bounds for
\emph{comparison-based} testers, \ie, testers that can only perform
comparisons of the function values at queried points, but cannot use
the values themselves.  Our main technical contribution is the
$\Omega(d\log d)$ lower bound for nonadaptive comparison-based testers
of unateness on hypercube domains.

Note that nonadaptivity must be critical in our lower bound construction, since
we obtained 
%
%
an 
%
$O(d)$-query adaptive tester for unateness.
%
%
%
%
Another challenge in proving the lower bound is the existence of a \emph{single, universal} nonadaptive
$O(d)$-tester for all
$\bb$-monotonicity properties, proven in~\cite{CDJS17}. In other words, there is a single
distribution on  $O(d)$ queries that defines a nonadaptive property tester for
$\bb$-monotonicity, regardless of $\bb$.
Since unateness
is the union of all $\bb$-monotonicity properties,
our construction of hard inputs must be able to fool such algorithms. In particular, if $\bb$ is fixed in advance, the algorithm from~\cite{CDJS17} will work, so it should be hard to learn $\bb$ for functions in the lower bound construction.
Once a tester finds a non-constant edge in each dimension, the problem reduces to testing $\bb$-monotonicity for a vector $\bb$ determined by the directions (increasing or decreasing) of the non-constant edges. That is, intuitively, most edges in our construction must be constant. This is one of the main technical challenges. The previous lower bound constructions for monotonicity testing \cite{BBM12,CS14} crucially used the fact that all edges in the hard functions were non-constant.

We briefly describe how we overcome the problems mentioned above.
By Yao's minimax principle, it suffices to construct two distributions,
%
$\Yes$ and $\No$, on unate and
%
on  %
far-from-unate functions,   %
%
respectively, %
that a deterministic nonadaptive tester cannot distinguish.
First, for some parameter $m$, we partition the hypercube
into $\totcube$ subcubes based 
%
%
on 
%
the first $\log_2\totcube$ most significant coordinates.
Both distributions, $\Yes$ and $\No$, sample a uniform $k$ from $[K]$, where $K=\Theta(\log d)$, and a set $R\subseteq[d]$ of cardinality $2^k$.
Furthermore, each subcube $j\in [\totcube]$ selects an ``action dimension'' $r_j\in R$ uniformly at random. For both distributions, in any particular subcube $j$, the function value is completely determined by the coordinates \emph{not in} $R$, and the random coordinate $r_j\in R$. Note that all the $i$-edges for $i\in (R\setminus \{r_j\})$ are constant.
Within the subcube, the function is a linear function
with exponentially increasing coefficients.
In the distribution $\Yes$, any two subcubes $j,j'$ with the same action dimension orient the edges in that dimension the same  way (both increasing or both decreasing), whereas in the distribution $\No$ each subcube decides on the orientation independently.
The former correlation maintains unateness while the 
latter independence creates distance to unateness.
We prove that to distinguish the distributions,
any comparison-based nonadaptive tester must find two distinct subcubes
with the same action dimension $r_j$ and, furthermore, make a specific
%
query (in both) that reveals the coefficient of $r_j$\,.
%
We show that, with $o(d\log d)$ queries, the probability of this event is negligible.

\section{Upper bounds}\label{sec:upper-bounds}
In this section, we prove parts 1-2 of \Thm{main-hg}, starting from the hypercube domain.

Recall the definition of $i$-edges and $i$-lines from \Sec{intro-tech} and what it means for an edge to be increasing, decreasing, and constant.

The starting point for our algorithms is the dimension reduction theorem from~\cite{CDJS17}. It bounds the distance of $f:[n]^d \to \R$ to monotonicity in terms of
the %
average distances of restrictions of $f$ to one-dimensional functions.
%
\begin{theorem}[Dimension Reduction, Theorem 1.8 in \cite{CDJS17}] \label{thm:dimred} Fix a bit vector ${\bf b}\in\{0,1\}^d$ and a function $f:[n]^d \to \R$ which is $\eps$-far from ${\bf b}$-monotonicity.
	For all $i\in[d]$, let $\mu_i$ be the average distance of $f_{|\ell}$ to $\bb_i$-monotonicity over all $i$-lines $\ell$.
	Then \[\sum_{i=1}^d \mu_i \geq \frac{\eps}{4}.\]
\end{theorem}
For the special case of the hypercube domains, $i$-lines become $i$-edges, and the average distance $\mu_i$ to $\bb_i$-monotonicity is the fraction of $i$-edges on which the function is not $\bb_i$-monotone.


\subsection{The nonadaptive tester over the hypercube}\label{sec:non-adap-ub}
We now describe \Alg{na-unate}, the nonadaptive tester for unateness over the hypercube 
domains. %

\begin{algorithm}
\caption{The Nonadaptive Unateness Tester over the Hypercube} \label{alg:na-unate}
\SetKwInOut{Input}{input}\SetKwInOut{Output}{output}
\SetKwFor{RepeatTimes}{repeat}{times}{end}
\SetKwFor{RepeatIterations}{repeat}{iterations}{end}
\Input{
a %
distance parameter $\eps \in (0,1/2)$; query access to a function $f:\{0,1\}^d \to \R$.}
\DontPrintSemicolon
\BlankLine
\nl \For{$r = 1$ to $\ceil{\log(16d/\eps)}$}{
\nl \RepeatTimes{$s_r = \ceil{\frac{16d \ln 4}{\eps\cdot 2^r}}$}{
\nl \label{step:sample-dim}Sample a dimension $i \in [d]$ uniformly at random. \;
\nl \label{step:reject-hc}Sample $3 \cdot 2^r$ $i$-edges uniformly and independently at random and {\bf reject} if there exist an increasing edge and a decreasing edge among the sampled edges.\;
}
}
\nl {\bf accept}\;
\end{algorithm}

It is evident that \Alg{na-unate} is a nonadaptive, one-sided error
tester. Furthermore, its query complexity is
$O\left(({d}/{\eps})\log({d}/{\eps})\right)$. It suffices to prove
the following.

\begin{lemma} \label{lem:na-unate-rej}
If $f$ is $\eps$-far from unate,
	\Alg{na-unate} rejects with probability at least $2/3$.
\end{lemma}

\begin{proof} Recall that $\alpha_i$ is the fraction of $i$-edges that are decreasing, $\beta_i$ is the fraction of $i$-edges that are increasing and $\mu_i = \min(\alpha_i,\beta_i)$.
	
	
	Define the $d$-dimensional bit vector ${\bf b}$ as follows: for each $i\in [d],$ let ${\bf b}_i = 0$
	if $\alpha_i < \beta_i$ and $\bb_i = 1$ otherwise.
	Then the average distance of $f$ to $\bb_i$-monotonicity over a random $i$-edge is precisely $\mu_i$.
	Since $f$ is $\eps$-far from being unate,
	$f$ is also $\eps$-far from being ${\bf b}$-monotone.
	By \Thm{dimred}, $\sum_{i \in [d]} \mu_i \geq \eps/4$. Hence, $\E_{i \in [d]}[\mu_i] \geq \eps/{(4d)}$.
	We now apply the work investment strategy due to Berman et al.~\cite{BerRY14} to get an upper bound on the probability that \Alg{na-unate} fails to reject.

\begin{theorem}[\cite{BerRY14}]\label{thm:wis}
For a random variable $X \in [0,1]$ with $\E[X] \geq \mu$, let $p_r = \Pr[X \geq 2^{-r}]$ and $\delta \in (0,1)$ be the desired error probability. Let $s_r = \frac{4 \ln 1/\delta}{\mu \cdot 2^r}$. Then
\[\prod\limits_{r=1}^{\ceil{\log (4/\mu)}} (1-p_r)^{s_r} \leq \delta. \]
\end{theorem}

Consider running \Alg{na-unate} on a function $f$ that is $\eps$-far from unate. Let $X = \mu_i$ where $i$ is sampled uniformly at random from $[d]$. Then $\E[X] \geq \eps/{(4d)}$. Applying the work investment strategy (\Thm{wis}) on $X$ with $\mu= \eps/{(4d)}$, we get that the probability that, in some iteration, \Step{sample-dim} samples a dimension $i$ such that $\mu_i\geq 2^{-r}$ is at least $1-\delta$.
We set $\delta = 1/4$.
Conditioned on sampling such a dimension, the probability that \Step{reject-hc} fails to obtain an increasing edge and a decreasing edge among its $3 \cdot 2^r$ samples is at most $2\left(1-2^{-r}\right)^{3 \cdot 2^r} \leq 2\eee^{-3} < 1/9$, as the fraction of both increasing and decreasing edges in the dimension is at least $2^{-r}$.
Hence, the probability that \Alg{na-unate} rejects $f$ is at least $({3}/{4}) \cdot ({8}/{9}) = {2}/{3}$,
which completes the proof of \Lem{na-unate-rej}.
\end{proof}

\subsection{The adaptive tester over the hypercube}\label{sec:adap-ub}
We now describe \Alg{adap-unate}, an adaptive tester for unateness over the hypercube domains with good expected query complexity. The final tester is obtained by repeating this tester and accepting if the number of queries exceeds a specified bound.

\begin{algorithm}
\caption{The Adaptive Unateness Tester over the Hypercube} \label{alg:adap-unate}
\SetKwInOut{Input}{input}\SetKwInOut{Output}{output}
\SetKwFor{RepeatTimes}{repeat}{times}{end}
\SetKwFor{RepeatIterations}{repeat}{iterations}{end}
\Input{
a %
distance parameter $\eps \in (0,1/2)$; query access to a function $f:\{0,1\}^d \to \R$.}
\DontPrintSemicolon
\BlankLine
\nl\label{step:repeat}\RepeatTimes{$\ceil{8/\eps}$}{
\nl\label{step:begin-adap-hc}\For{$i = 1$ to $d$}{
\nl\label{step:sample-s}Sample an $i$-edge $e_i$ uniformly at random.\;
\nl\label{step:end-build-s}\If{$e_i$ is non-constant (\ie, increasing or decreasing)}{
\nl\label{step:stage-2-edge-sample} Sample $i$-edges uniformly at random until we obtain a non-constant edge $e_i'$.\;
\nl\label{step:end-adap-hc} {\bf reject} if one of the edges $e_i$ and $e_i'$ is increasing and the other is decreasing.\;}
}}
\nl {\bf accept}
\end{algorithm}

\begin{claim} \label{clm:time}
The expected number of queries made by
\Alg{adap-unate} is $O(d/\eps)$.
\end{claim}

\begin{proof}
Consider one iteration of the {\bf repeat}-loop in \Step{repeat}.
We prove that the expected number of queries in this iteration is $4d$.
The total number of queries in \Step{sample-s} is $2d$, as 2 points per dimension are queried.
Let $E_i$ be the event that edge $e_i$ is non-constant and $T_i$ be the random variable for the number of $i$-edges sampled in \Step{stage-2-edge-sample}. Then
\[
  \E[T_i] = \frac 1{\alpha_i+\beta_i}=\frac 1{\Pr[E_i]}\,.
\]
Therefore, the expected number of all edges sampled in
\Step{stage-2-edge-sample} is
\[
  \sum_{i=1}^d \Pr[E_i]\cdot\E[T_i]=\sum_{i=1}^d \Pr[E_i] \cdot \frac
  1{\Pr[E_i]} = d\,.
\]
Hence, the expected number of queries in \Step{stage-2-edge-sample} is
$2d$, and the total expected number of queries in one iteration is
$4d$. Since there are $\ceil{8/\eps}$ iterations in \Step{repeat}, the
expected number of queries in \Alg{adap-unate} is $O(d/\eps)$.
\end{proof}

\begin{claim} \label{clm:rej}
If $f$ is $\eps$-far from unate, \Alg{adap-unate} accepts with probability at most $1/6$.
\end{claim}

\begin{proof}
First, we bound the probability that a violation of unateness is detected in some dimension $i \in [d]$ in one iteration of the {\bf repeat}-loop in \Step{repeat}.
Consider the probability of finding a %
decreasing $i$-edge in \Step{sample-s} and of finding an %
increasing $i$-edge in \Step{stage-2-edge-sample}.
The former is exactly $\alpha_i,$ and the latter is ${\beta_i}/{(\alpha_i+\beta_i)}$.
Similarly, the probability of finding an increasing $i$-edge in \Step{sample-s} and of finding a decreasing $i$-edge in \Step{stage-2-edge-sample} is $\beta_i$ and ${\alpha_i}/{(\alpha_i+\beta_i)}$, 
respectively.  %
Therefore, the probability we detect a violation from dimension $i$ is
\[2 \cdot \frac{\alpha_i \beta_i}{\alpha_i+\beta_i} \geq \min(\alpha_i,\beta_i) = \mu_i.\]
The probability that we fail to detect a violation in any of the $d$ dimensions in one iteration is at most
\[\prod_{i =1}^d (1-\mu_i) \leq \exp\big(-\sum_{i = 1}^d \mu_i\big) \leq \eee^{-\eps/4},\]
where the last inequality follows from \Thm{dimred} (Dimension Reduction).
Thus, the probability that \Alg{adap-unate} fails to reject in all iterations of \Step{repeat} is at most
$(\exp(-\eps/4))^{(8/\eps)} = \eee^{-2} < 1/6$.
\end{proof}

\begin{proof}[{Proof of \Thm{main-hg}, Part 2} (for the special case of the hypercube domain)]
We run \Alg{adap-unate}, aborting and accepting if the number of queries exceeds the expectation by a factor of 6. By Markov's inequality, the probability of aborting is at most $1/6$.
By \Clm{rej}, if $f$ is $\eps$-far from unate, \Alg{adap-unate} accepts with probability at most $1/6$. The theorem follows by a union bound.
\end{proof}

\subsection{Extension to hypergrids}\label{sec:hypergrids}
We start by establishing terminology for lines and pairs.
Consider a function $f:[n]^d\to\R$.
Recall the definition of $i$-lines from \Sec{intro-tech}.
A pair of points that differ only in coordinate $i$ is called an $i$-pair.
An $i$-pair $\{x,y\}$ with $x_i < y_i$ is called \emph{increasing} if $f(x) < f(y)$, \emph{decreasing} if $f(x) > f(y),$ and \emph{constant} if $f(x) = f(y)$. Recall that we call a function $h:[n] \to \R$ \emph{monotone} if $h(x)\leq h(y)$ for all $x<y$ and 
%
%
\emph{antitone} %
if $h(x)\geq h(y)$ for all $x < y$.

Recall that a proximity-oblivious tester (POT) for a property $\cP$ rejects every input with probability proportional to the distance from the input to $\cP$.
As discussed in \Sec{intro-tech}, there are several 1-sided error, nonadaptive POTs for monotonicity on the line that we can use to extend Algorithms~\ref{alg:na-unate} and~\ref{alg:adap-unate} to work on hypergrids. One of them is the \emph{tree tester}, designed by Ergun et al.~\cite{EKKRV00}, that simply picks a point $x\in[n]$ uniformly at random, queries all points visited in a binary search for $x$, and rejects iff $x$ forms a decreasing pair with one of the queried points. By symmetry, when this algorithm is modified to reject iff it finds an increasing pair, it tests 
%
%
antitonicity. %
The tree tester (and any other 1-sided error POT for sortedness) can be easily modified to return whether it found any increasing/decreasing edges by including $\uparrow$ / $\downarrow$ in its output.

\begin{lemma}[\cite{EKKRV00,BGJRW12,CS13,CDJS17}] \label{lem:line}
%
There exists a nonadaptive algorithm $\mpot$ that gets query access to a function $h:[n] \to \R$, makes $O(\log n)$ queries, and returns a set  $\dir\subseteq\{\uparrow,\downarrow\}$.
%
%
%
\begin{itemize}
\item If $h$ is monotone, then $\downarrow$ is not in $\dir$. The probability that $\dir$ contains $\downarrow$ is at least the distance from $h$ to monotonicity.
\item Similarly, if $h$ is 
%
%
antitone, %
then $\uparrow$ is not in $\dir$. The probability that $\dir$ contains $\uparrow$ is at least the distance from $h$ to
%
%
antitonicity. %
\end{itemize}
%
\end{lemma}

We now describe \Alg{adap-unate-hg}, an adaptive tester for unateness over the hypergrid domains with good expected query complexity. As in the case of the hypercube 
domains, %
the final tester is obtained by repeating this tester and accepting if the number of queries exceeds a specified bound.


\begin{algorithm}
\caption{The Adaptive Unateness Tester over the Hypergrid} \label{alg:adap-unate-hg}
\SetKwInOut{Input}{input}\SetKwInOut{Output}{output}
\SetKwFor{RepeatTimes}{repeat}{times}{end}
\SetKwFor{RepeatIterations}{repeat}{iterations}{end}
\Input{
a %
distance parameter $\eps \in (0,1/2)$; query access to a function $f:[n]^d \to \R$.}
\DontPrintSemicolon
\BlankLine
\nl \RepeatTimes{$\ceil{8/\eps}$}{
\nl\label{step:begin-adap-hg}\For{$i = 1$ to $d$}{
\nl Sample an $i$-line $\ell_i$ uniformly at random.\;
\nl\label{step:test-i-line}Let $\dir_i$ be the output of $\mpot$ (from \Lem{line}) on $f_{|\ell_i}$. \;
\nl \If{$\dir_i \ne \emptyset$}{
%
%
%
\nl \Repeat{$\dir_i'\neq\emptyset$}{
\nl\label{step:verify-i-line} Sample an $i$-line $\ell_i'$ uniformly at random and let $\dir'_i$ be the output of $\mpot$  on $f_{|\ell'_i}$.\;}
\nl  If $\dir_i \cup \dir_i' = \{\uparrow, \downarrow \}$, {\bf reject}.\;
}}}
\nl {\bf accept}
\end{algorithm}


\begin{proof}[Proof of \Thm{main-hg}, Part 2]
First, we explain how \Lem{line} and \Thm{dimred} are used in the analysis of the adaptive tester.
For a dimension $i\in[d]$, let $\alpha_i$ and $\beta_i$ denote the average distance of $f_{|\ell}$ to monotonicity and
%
%
antitonicity, %
respectively, %
over all $i$-lines $\ell$.
Then $\mu_i := \min(\alpha_i,\beta_i)$ is the average fraction of points per $i$-line that needs to change to make $f$ unate.
Define the $\bb$-vector with $\bb_i = 0$ if $\alpha_i < \beta_i$,
and $\bb_i = 1$ otherwise.
By \Thm{dimred},
if $f$ is $\eps$-far from unate, and thus $\eps$-far from $\bb$-monotone,
then $\sum_{i=1}^d \mu_i \geq \eps/4$.
%
By \Lem{line},
the probability that the output of $\mpot$ on a uniformly random  
%
$i$-line $f_{|\ell_i}$ contains $\downarrow$
%
is at least $\alpha_i$\,,
and the probability that it contains $\uparrow$ is at least $\beta_i$\,.
%
%
%
The rest of the analysis of \Alg{adap-unate-hg} is similar to that in the 
hypercube case.

In the proof of \Clm{time}, the expected number of edges sampled in one iteration of \Alg{adap-unate} is $2d$. Similarly, the expected number of lines sampled in one iteration of \Alg{adap-unate-hg} is $2d$. Since $\mpot$ makes $O(\log n)$ queries, the overall expected number of queries is $O((d\log n)/\eps)$.

The proof of \Clm{rej} carries over almost word for word. Fix a dimension $i$. The probability that $\downarrow \; \in \dir_i$ in \Step{test-i-line} is at least $\alpha_i$.
The probability that $\uparrow \; \in \dir_i'$ in \Step{verify-i-line} is at least ${\beta_i}/{(\alpha_i + \beta_i)}$.
The rest of the calculation is identical to that of the proof of \Clm{rej}.

Finally, we run \Alg{adap-unate-hg} and abort and accept if the number of queries exceeds the expectation by a factor of 6. As in the proof of \Thm{main-hg}, Part 2, the resulting algorithm always accepts unate functions and accepts functions that are $\eps$-far from unate with probability at most $1/3$, completing the~proof.
\end{proof}


%

Next we show how to modify any POT for sortedness in a black-box manner to obtain a POT for unateness on the line with the same guarantees.

\begin{algorithm}
\caption{The POT for Unateness over the Line}\label{alg:unateness-line}
\SetKwInOut{Input}{input}\SetKwInOut{Output}{output}
\SetKwFor{RepeatTimes}{repeat}{times}{end}
\SetKwFor{RepeatIterations}{repeat}{iterations}{end}
\Input{
query %
access to a function $h:[n] \to \R$.}
\DontPrintSemicolon
\BlankLine
\nl Run algorithm $\mpot$ on $h$ and let $\dir$ be its output. \;
\nl Pick $x \in [n]$ uniformly at random and query $h$ on points $1,x,$ and $n$. \label{step:pick-x} \;
\nl If any pair in $\{1,x,n\}$ is increasing, update $\dir\gets\dir\cup\{\uparrow\}$. \;
\nl If any pair in $\{1,x,n\}$ is decreasing, update $\dir\gets\dir\cup\{\downarrow\}$. \;
\nl {\bf Reject} if $\dir=\{\uparrow,\downarrow\}$; otherwise, {\bf accept}.
\end{algorithm}

\begin{lemma}\label{lem:line-unateness}
If $h:[n] \to \R$ is $\eps$-far from unate, then \Alg{unateness-line} rejects with probability at least~$\eps$.
\end{lemma}
\begin{proof}
First, consider the case when $h(1)=h(n)$. Since $h$ is $\eps$-far from unate, it is also $\eps$-far from being a constant function equal to $h(1)$ on all points in $[n]$. Therefore, with probability at least $\eps$, point $x$ chosen in \Step{pick-x} of the algorithm satisfies $h(x)\neq h(1)$. If this holds, then one of the pairs $\{1,x\}$ and $\{x,n\}$ is increasing and the other is decreasing. So, \Alg{unateness-line} indeed rejects with probability at least $\eps$.

Next, consider the case when $h(1)<h(n)$. Then $\{1,n\}$ is an increasing pair. Since $h$ is $\eps$-far from monotone, by \Lem{line}, the output of $\mpot$ contains $\downarrow$ with probability at least $\eps$. So, again \Alg{unateness-line} rejects with probability at least $\eps$.

Finally, the case when $h(1)>h(n)$ is symmetric to the case when $h(1)<h(n)$.
\end{proof}




%

\begin{algorithm}[h!]
	\caption{The Nonadaptive Unateness Tester over the Hypergrid} \label{alg:na-unate-hg}
	\SetKwInOut{Input}{input}
	\SetKwFor{RepeatTimes}{repeat}{times}{end}
	\SetKwFor{RepeatIterations}{repeat}{iterations}{end}
	\Input{distance parameter $\eps \in (0,1/2)$; query access to a function $f:[n]^d \to \R$.}
	\DontPrintSemicolon
	\BlankLine
	
	\nl \label{step:mono}\RepeatTimes{$\ceil{\frac{24\ln 3}\eps}$}{
		\nl \label{step:mono-start}\For{$i=1$ to $d$}{
			\nl Sample an $i$-line $\ell$ uniformly at random. \;
			\nl \label{step:log}{\bf Reject} if \Alg{unateness-line} rejects on input $f_{|\ell}$.
		}
	}
	\nl \label{step:cube}\For{$r=1$ to $\ceil{\log(96d/\eps)}$}{
		\nl \RepeatTimes{$s_r = \ceil{\frac{96d \ln 4}{\eps \cdot 2^r}}$}{
			\nl \label{step:samp-hg1}Sample a dimension $i \in [d]$ uniformly at random.\;
			\nl \label{step:tree-tester1}Sample $3 \cdot 2^r$ $i$-pairs uniformly and independently at random.\;
			\nl \label{step:rej-hg1}If we find an increasing and a decreasing pair among the sampled pairs, {\bf reject}.\;
		}
	}
	\nl {\bf accept}			
\end{algorithm}



Our nonadaptive hypergrid tester is stated in \Alg{na-unate-hg}.
A crucial part of its analysis is \Lem{dychotomy} that demonstrates that every function on the line that is far from monotone is also far from unate or a large fraction of pairs are decreasing with respect to that function.

\begin{lemma}\label{lem:dychotomy}
%
%
%
%
Consider a function $h:[n]\to \R$ which is $\eps$-far from
monotone.
%
%
%
%
%
At least one of the following holds:
\begin{enumerate}
\item
%
The function $h$  %
is $\frac \eps 2$-far from unate.
%
%
\item $\Pr_{u,v \in [n]}\Big[\text{pair } \{u,v\} \text{ is decreasing
}\Big] \geq \frac\eps 4$.
%
%
%
\end{enumerate}
\end{lemma}
\begin{proof}
%
%
Suppose
%
Item~1 does not hold, that is, there exists a set $G\subseteq [n]$ of size greater than $n(1-\eps/2)$ on which $h_{| G}$ is unate and therefore 
%
%
antitone.
%
%
(Unate means monotone or 
%
%
antitone, %
but $h_{|G}$ cannot be monotone because $h$ is $\eps$-far from monotone.) Observe that if $u,v\in G$ and $h(u)\neq h(v)$ then $\{u,v\}$ is a decreasing pair because $h_{|G}$ is 
%
%
antitone. 
%
For each point $u\in G$, let $D_u\subseteq G$ be the set of points on which $h$ differs from $h(u)$. Since $h$ is $\eps$-far from monotone and, consequently, $\eps$-far from constant, $|D_u|\geq\eps/2$ for all $u\in G$. Thus,
\[
\Pr_{u,v \in [n]}[\text{pair } \{u,v\} \text{ is decreasing}]
\geq \Pr[u\in G \text{ and } v\in D_u]\geq \Big(1-\frac\eps 2\Big)\cdot \frac \eps 2\geq \frac \eps 4.
\]
The last inequality follows because $\eps\leq 1$. We proved that if Item~1 does not hold then Item~2 must hold.
\end{proof}

\begin{theorem} \label{thm:non-adap-hyper} If $f:[n]^d \to \R$ is $\eps$-far
	from unate, then \Alg{na-unate-hg} rejects with probability at least $2/3$.
\end{theorem}

\begin{proof} Suppose $f:[n]^d \to \R$ is $\eps$-far from unate.
We will show that when \Alg{na-unate-hg} runs on $f$, one of the two loops
%
%
(in Steps~\ref{step:mono} and~\ref{step:cube})   %
rejects with probability at least 2/3.

For every line $\ell$, we define the following quantities.
\begin{itemize}
\item $\alpha_\ell$: the distance of $f_{|\ell}$ to monotonicity.
\item $\beta_\ell$: the distance of $f_{|\ell}$ to
%
%
antitonicity. %
\item $\mu_\ell=\min(\alpha_\ell,\beta_\ell)$: the distance of $f_{|\ell}$ to unateness.
%
\item $\alpha'_\ell$: the probability that a uniformly random pair in $\ell$ is decreasing.
\item $\beta'_\ell$: the probability that a uniformly random pair in $\ell$ is increasing.
%
\end{itemize}

%
%
By \Lem{dychotomy} and symmetry, for every line $\ell$,
%
\begin{align}
	\mu_\ell + 2\alpha'_\ell \geq \frac{\alpha_\ell}2 \text{ and } \mu_\ell + 2\beta'_\ell \geq \frac{\beta_\ell}2. \label{eqn:dychotomy}
\end{align}

Let $L_i$ be the set of $i$-lines. By
\Thm{dimred},
\begin{align}
\frac{1}{n^{d-1}} \sum_{i=1}^d \min\left(\sum_{\ell \in L_i} \alpha_\ell, \sum_{\ell \in L_i} \beta_\ell \right) \geq \frac{\eps}{4}. \label{eqn:dim-reduction}
\end{align}

%
%

Define $u=\sum_{i=1}^d u_i,$ where
\[
  u_i= \frac{1}{n^{d-1}}\sum_{\ell \in L_i} \mu_\ell
\]
is the average distance to unateness in dimension $i$. (In general,
$u_i$ is different from the quantity $\mu_i$ used in the analysis of
the adaptive tester. Recall that $\mu_i$ was used to denote the
minimum of the average distance to monotonicity and the average
distance to 
%
%
antitonicity %
in dimension $i$. We are using different
notation to avoid confusion.)  Let $\alpha'_i$ to be the average of
$\alpha'_\ell$ for all lines $\ell\in L_i$. Define $\beta'_i$
analogously, and let $\mu'_i=\min(\alpha'_i,\beta'_i)$. Let random
variable $X$ be $\mu'_i$ where $i$ is sampled uniformly at random from
$[d]$.  Then
\begin{align*}
u+2d\cdot\E[X]
&=	\frac{1}{n^{d-1}}\sum_{i=1}^d  \left[\sum_{\ell \in L_i} \mu_\ell + 2\cdot\min\left(\sum_{\ell \in L_i} \alpha'_\ell, \sum_{\ell \in L_i} \beta'_\ell \right)\right]  \\
	&= \frac{1}{n^{d-1}}\sum_{i=1}^d  \min\left(\sum_{\ell \in L_i} (\mu_\ell + 2\alpha'_\ell), \sum_{\ell \in L_i} (\mu_\ell + 2\beta'_\ell) \right)\\
&\geq \frac{1}{n^{d-1}}\sum_{i=1}^d  \min\left(\sum_{\ell \in L_i} \frac{\alpha_\ell}2, \sum_{\ell \in L_i} \frac{\beta_\ell}2 \right)
\geq \frac{\eps}8.
	\end{align*}
The last two inequalities follow from Equations~\eqref{eqn:dychotomy}
and~\eqref{eqn:dim-reduction}, respectively.
%
%
We conclude that at least one of $u$ and $\E[X]$ has to be
large, %
specifically, $u \geq \eps/{24}$ or $\E[X]\geq \eps/{(24d)}$.
	
{\bf Case 1:} $u\geq \eps/{24}$. Consider the first loop of \Alg{na-unate-hg} (in \Step{mono}). By \Lem{line-unateness}, the probability that a uniformly random $i$-line is rejected by \Alg{unateness-line} is at least $u_i$. The probability that one iteration of the loop in \Step{mono} fails to reject is
\[\prod_{i=1}^d (1-u_i)
\leq \prod_{i=1}^d \exp(-u_i)
=\exp(-\sum_{i=1}^d u_i)
=\exp(-u)\leq \exp \big(-\frac \eps {24}\big).
\]
Thus, all iterations of the loop fail to reject with probability at most $(\exp(-\eps/24))^{(24\ln 3)/\eps}=1/3$.

{\bf Case 2:} $\E[X]\geq \eps/{(24d)}$.
Applying the work investment strategy (\Thm{wis}) on $X$ with $\mu=\eps/{(24d)}$ and using a calculation analogous to that in the proof of \Lem{na-unate-rej},
the probability that \Step{rej-hg1} rejects in some iteration is at least
$2/3$.
\end{proof}


\section{The lower bound for nonadaptive testers over the hypercube} \label{sec:lb}
In this section, we prove \Thm{non-adap-lb-1}, which gives a lower bound for nonadaptive unateness testers for functions over the hypercube.

Fischer~\cite{Fis04} showed that in order to prove lower bounds for a general class of properties on the line domain, it is sufficient to consider a special class of testers called \emph{comparison-based testers}. The properties he looked at are called \emph{order-based properties} (see \Def{obp}), and they include monotonicity and unateness.
A tester is \emph{comparison-based} if it bases its decisions only on the \emph{order} of the function values at the points it queried,  and not on the values themselves.
Chakrabarty and Seshadhri~\cite{CS14} extended Fischer's proof to monotonicity on any partially-ordered domain for the case when all function values are distinct. As we show in \Sec{reduction-cbt} below, Chakrabarty and Seshadhri's proof goes through for all order-based properties on partially-ordered domains. Moreover, the assumption of distinctness for function values can be removed. We include this proof for completeness, filling in the details needed to generalize the original proof.

Our main technical contribution is the construction of a distribution of functions $f:\{0,1\}^d \to \N$ on which every nonadaptive comparison-based tester must query $\Omega(d \log d)$ points to determine whether the sampled function is unate or far from unate.
We describe this construction in \Sec{hard-dist} and show its correctness in Sections~\ref{sec:sgn-fn}-\ref{sec:bad}.

\subsection{Reduction to comparison-based testers}\label{sec:reduction-cbt}
In this section, we prove that if there exists a tester for an order-based property ${\cal P}$ of functions over a partially-ordered domain, then there exists a comparison-based tester for ${\cal P}$ with the same query complexity. This is stated in \Thm{CS14}.
Before stating the theorem, we introduce some definitions.

\begin{definition} \label{def:t-eps-del-tester}
A $(t,\eps,\delta)$-tester for a property is a 2-sided error tester with distance parameter $\eps$ that makes at most $t$ queries and errs with probability at most $\delta$.
\end{definition}

\begin{definition}[Order-based property]\label{def:obp}
For an arbitrary partial order $D$ and an arbitrary total order $R$, a property $\cP$ of functions $f:D \to R$ is \emph{order-based} if,
for all strictly increasing maps $\phi: R \to R$ and all functions $f$, we have $\dist(f,\cP) = \dist(\phi\circ f, \cP)$.
\end{definition}

In particular, unateness is an order-based property.
The following theorem is an extension of Theorem~5 in~\cite{Fis04} and Theorem~2.1 in~\cite{CS14}.
Specifically, Theorem~2.1 in~\cite{CS14} was proved with the assumption that the function values are distinct. We generalize
%
the theorem
by removing this assumption.

\begin{theorem}[Generalization of~\cite{Fis04,CS14}]\label{thm:CS14}
Let $\cP$ be an order-based property of functions $f:D \to \N$ over a finite domain $D$. Suppose there exists a $(t,\eps,\delta)$-tester for $\cP$.
Then there exists a comparison-based $(t,\eps,2\delta)$-tester
for $\cP$.
\end{theorem}

The rest of this section is devoted to proving \Thm{CS14}.
Our proof closely follows the proof of Theorem~2.1 in~\cite{CS14}.
The proof has two parts. The first part describes a reduction from a tester to a \emph{discretized tester}, and the second part describes a reduction from a discretized tester to a comparison-based tester.

Let $\cP$ be a property of functions $f:D \to R$ for an arbitrary partial order $D$ and an arbitrary total order $R \subseteq \N$. Let $\cT$ be a $(t,\eps,\delta)$-tester for $\cP$. First, we define a family of probability functions that completely characterizes $\cT$. Fix some $s \in [t]$. Consider the point in time in an execution of the tester $\cT$ on some input function $f$, where exactly $s$ queries have been made.
Suppose these queries are $x_1, x_2, \ldots, x_s \in D$ and the corresponding answers are $a_1=f(x_1), a_2=f(x_2), \ldots, a_s=f(x_s)$.
Let the \emph{query vector} $X$ be $(x_1, \ldots, x_s)$ and the \emph{answer vector} $A$ be $(a_1, \ldots, a_s)$. The next action of the algorithm is either choosing the ${(s+1)}^\text{st}$ query from $D$ or outputting \emph{accept} or \emph{reject}. For each action $y \in D \cup \{\texttt{accept}, \texttt{reject}\}$, let $p_X^y(A)$ denote the probability that $\cT$ chooses action $y$ after making queries $X$ and receiving answers $A$. Since $p_X^y(A)$ is a probability distribution,
\[\forall s < t, \forall X \in D^s, \forall A \in R^s, \sum_{y \in D \cup \{\texttt{accept}, \texttt{reject} \}} p_X^y(A) = 1.\]
%
Furthermore, the tester cannot make more than $t$ queries,
and so the action $(t+1)$ must be either \texttt{accept} or \texttt{reject}. Formally,
\[\forall X \in D^t, \forall A \in R^t, \sum_{y \in \{\texttt{accept}, \texttt{reject} \}} p_X^y(A) = 1.\]
%
If a tester decides to accept or reject before making $t$ queries, \ie, $p_X^{y'}(A) = 1$ for some $X = (x_1, \ldots, x_s), A=(a_1, \ldots, a_s)$, where $s < t$ and $y' \in \{\texttt{accept, reject}\}$, then
we fill in the values for $p$, so that the action of the tester is defined to be the same $y'$ for all values until $t+1$. Specifically,
 we set $p_{X'}^{y'}(A') = 1$
for all $X' \in D^{s'}, A' \in R^{s'}$ where $s < s' \leq t$ and the first $s$ queries (in $X'$) and their corresponding answers (in $A'$) are $x_1, \ldots, x_s$ and 
$a_1, \ldots, a_s$\,, %
respectively.  %

\begin{definition} [Discretized tester] \label{def:disc}
A tester $\cT$ is \emph{discretized} if all $p_X^y(A)$-values associated with $\cT$ come from the range $\left\{ {i}/{K}: i\in \{0,1,\ldots,K\} \right\}$ for some positive integer $K$.
\end{definition}

Chakrabarty and Seshadhri~\cite{CS14} proved that if there exists a $(t,\eps,\delta)$-monotonicity tester $\cT$ for functions $f:D \to \N$, then there exists a discretized $(t,\eps,2\delta)$-monotonicity tester $\cT'$ for the same class of functions.
Both the statement and the proof in~\cite{CS14} hold not only for testers of monotonicity, but for testers of all properties of functions $f:D \to R$.

\begin{lemma}[Implicit in~{\cite[Lemma~2.2]{CS14}}] \label{lem:disc-test}
Suppose there exists a $(t,\eps,\delta)$-tester $\cT$ for a property $\cP$ of functions $f:D \to R$.
Then there exists a $(t,\eps,2\delta)$-discretized tester $\cT'$ for $\cP$.
\end{lemma}

This completes the first part of the proof.

Next, we show how to transform a discretized tester into a comparison-based tester.
Intuitively, a tester is comparison-based if each action of the tester depends only on the ordering of the answers to the previous queries, not on the values themselves.
We define a family of probability functions $q$ in order to characterize comparison-based testers.
The $q$-functions are defined in terms of $p$-functions, but, in their definition, we decouple the set of values that were received as answers from their positions in the answer vector. Specifically, the set of values that were received as answers (\ie, information irrelevant for
a %
comparison-based tester) is given as the argument to the $q$-functions. All the remaining information is given as subscripts and superscripts.
Let $V$ represent the set $\{a_1, \ldots, a_s\}$ of answer values (without duplicates).
Let $r$ be the number of (distinct) values in $V$. Note that $r \leq s$.
Suppose $V$ is $\{v_1, v_2, \ldots, v_r\},$ where $v_1, \ldots, v_r \in R$ 
and $v_1 < v_2 < \ldots < v_r$\,. %
Let $\rho$ be the map from positions of values in the answer vector to their corresponding indices in $V$, that is, $\rho:[s] \to [r]$. Observe that $\rho$ is surjective.
The $q$-functions are defined as follows:
\[q_{X,\rho}^y(V) = p_X^y((v_{\rho(1)}, v_{\rho(2)}, \ldots,v_{\rho(s)})).\]


For any set $S$, let $S^{(r)}$ denote the set of all subsets of $S$ of size $r$.

\begin{definition} [Comparison-based tester]\label{def:cbt}
A tester $\cT$ for an order-based property $\cP$ is \emph{comparison-based} for functions $f:D \to R$, if for all $r,s$ satisfying $ r \leq s \leq t$, and all $X \in D^s$, $y \in D \cup \{\texttt{accept}, \texttt{reject}\}$ and surjections $\rho:[s] \to [r]$, the function $q_{X,\rho}^y$ is constant on $R^{(r)}$. That is, for all $V, V' \in R^{(r)}$, we have $q_{X,\rho}^y(V) = q_{X,\rho}^y(V')$.
\end{definition}

To complete the proof of \Thm{CS14}, we show that if there exists a discretized tester $\cT$ for an order-based property $\cP$ over the functions $f:D \to \N$, then there exists an infinite set $R \subseteq \N$ such that, for functions $f:D \to R$, the tester $\cT$ is comparison-based. Specifically, the $q$-functions that describe the tester do not depend on $V$, the specific set of answer values, as long as $V\subset R$. At the end of the proof, we construct a new comparison-based tester that modifies the input function so that its range is $R$ (without changing the distance of the function to the property $\cP$) and runs $\cT$ on the modified function.
The existence of the infinite set $R$ is proved using Ramsey theory arguments.

We introduce some Ramsey theory terminology.
Consider a positive integer $C$, where $[C]$ represents a set of colors.
For any positive integer $i$, a \emph{finite coloring} of $\N^{(i)}$ is a function $\texttt{col}_i:\N^{(i)} \to [C]$. That is, this function assigns one of $C$ colors to each subset of $\N$ with $i$ elements. 
An infinite set $R \subseteq \N$ is \emph{monochromatic} with respect to $\texttt{col}_i$ if for all $i$-element subsets $V, V' \in R^{(i)}$, the color $\texttt{col}_i(V) = \texttt{col}_i(V')$. In other words, each $i$-element subset of $R$ is colored with the same color by the coloring function $\texttt{col}_i$.
 A \emph{$k$-wise finite coloring of $\N$} is a collection of $k$-colorings $\texttt{col}_1, \texttt{col}_2, \ldots, \texttt{col}_k$.
Note that each coloring function $\texttt{col}_1, \ldots, \texttt{col}_k$ is defined over subsets of a different size, and together they assign a color to each subset with at most $k$ elements.
An infinite subset $R \subseteq \N$ is \emph{$k$-wise monochromatic}
with respect to $\texttt{col}_1, \ldots, \texttt{col}_k$ if $R$ is monochromatic with respect to all $\texttt{col}_i$ for $i \in [k]$. That is, all subsets of $R$ of the same size get assigned the same color by the coloring functions.

We use the following variant of Ramsey's theorem which was also used in~\cite{Fis04,CS14}.

\begin{theorem}[Theorem~2.3 in~\cite{CS14}]\label{thm:ramsey}
For any $k$-wise finite coloring of $\N$, there exists an infinite $k$-wise monochromatic subset $R \subseteq \N$.
\end{theorem}


\begin{proof}[Proof of \Thm{CS14}]
Suppose there exists a $(t,\eps,\delta)$-tester for property $\cP$ of functions $f:D \to \N$. By \Lem{disc-test}, there exists a $(t,\eps,2\delta)$-discretized tester $\cT$ for $\cP$.
Consider the family of $q$-functions that characterizes $\cT$.

The main idea of the proof is to view the behavior of the tester $\cT$ on each possible subset of answer values $V$ (with at most $t$ elements) as the color of $V$. More precisely, the color of $V$ is the corresponding vector of all $q$-functions evaluated at $V$. If two sets $V$ and $V'$ are mapped to the same color, it means, intuitively, that the tester $\cT$ is ignoring whether the specific answer values are $V$ or $V'$. We want to show that there is a large subset $R$ of values, such that the tester $\cT$ ignores the answer values, as long as they come from $R$.
%
Since $\cT$ is discretized, the set of colors is finite, and we are able to apply \Thm{ramsey} to get an infinite subset $R$  of $\N$ which is $t$-wise monochromatic. That means that $\cT$ ignores $V$, as long as $V\subset R$. That is, $\cT$ is already comparison-based on $R$. At the end, we use the fact that $\cP$ is an order-based property to get a comparison-based tester for the whole range.

We define a $t$-wise finite coloring of $\N$. For each $r \in [t]$ and $V \in \N^{(r)}$, the color $\texttt{col}_r(V)$ is defined as a vector of probability values $q_{X,\rho}^y(V)$.
The vector is indexed by $(y, X, \rho)$ for each $y \in D \cup \{\texttt{accept}, \texttt{reject}\},  X \in D^s$ for an integer $s$ satisfying $r \leq s \leq t$, and a surjection
$\rho:[s] \to [r]$.
The value at the index $(y,X,\rho)$ in $\texttt{col}_r(V)$ is equal to $q_{X, \rho}^y (V)$.
Note that, there are finitely many possible values for $y$ and $X$, and surjections $\rho$.
%
%
So, the dimension of the vector $\texttt{col}_r(V)$ is finite. Furthermore, since the tester is discretized, the number of different values that the $q$-functions take is also finite. Hence, the range of $\texttt{col}_r$ is finite.
Now, we have a $t$-wise finite coloring $\texttt{col}_1, \ldots, \texttt{col}_t$ of $\N$.
By \Thm{ramsey}, there exists an infinite $t$-wise monochromatic set $R \subseteq \N$.
Thus, for each $r \in [t]$ and $V, V' \in R^{(r)}$, we have $\texttt{col}_r(V) = \texttt{col}_r(V')$, implying that $q_{X,\rho}^y (V) = q_{X,\rho}^y (V')$ for all $y, X,$
and %
$\rho$. Thus, $\cT$ is comparison-based for functions $f:D \to R$.

Finally, we construct a comparison-based tester $\cT'$ for the whole range $\N$. Consider a strictly monotone increasing map $\phi: \N \to R$. Given any function $f: D \to \N$, consider $\phi \circ f : D \to R$.
Define an algorithm $\cT'$, which on input $f$, runs $\cT$ on $\phi \circ f$.
Since $\cP$ is order-based, $\text{dist}(f, \cP) = \text{dist}(\phi \circ f, \cP)$.
Hence, $\cT'$ is a $(t,\eps, 2\delta)$-tester for $\cP$.
Moreover, since the tester $\cT'$ just runs $\cT$ on $\phi \circ f: D \to R$, and $\cT$ is comparison-based for $\phi \circ f$, the tester $\cT'$ is also comparison-based.
\end{proof}


\subsection{The hard distributions} \label{sec:hard-dist}
Our main lower bound theorem is stated next. Together with \Thm{CS14}, it implies \Thm{non-adap-lb-1}.

\begin{theorem}[Main]\label{thm:non-adap-lb}
Any nonadaptive comparison-based unateness tester of functions $f:\{0,1\}^d\to \N$ that works with distance parameter $\eps=1/8$ and errs with probability at most 1/6 must make $\Omega(d\log d)$ queries.
\end{theorem}

The proof of \Thm{non-adap-lb} is presented in Sections~\ref{sec:hard-dist}-\ref{sec:bad} and forms the core technical content of this work.
We will use Yao's minimax principle~\cite{Yao77}, specifically, the version stated in~\cite[Claim~5]{RS06}. It asserts that to prove a lower bound for a randomized tester, it is sufficient to give two distributions,
%
%
%
%
one on positive and one on negative instances,
%
that every deterministic tester fails to distinguish 
with high probability.
%
%
Next, we define two input distributions.

It may be useful for the reader to recall the  sketch of the main 
ideas given in \Sec{intro-tech}.
%
Without loss of generality,\footnote
{If the number of dimensions is $d''$ where there is no positive integer $d$ such that $d'' = d + \log_2 d$, then consider the largest integer $d'$ such that $d' < d''$ and $d' = d + \log_2 d$. We can construct any function over $\{0,1\}^{d'}$ and extend it to $\{0,1\}^{d''}$ by adding $d'' - d'$ dummy variables.
%
Furthermore, $d' = \Theta(d'')$, and hence, $\Omega(d'' \log d'') = \Omega(d' \log d')$.}
let $d$ be an even power of $2$ and $d' := d+\log_2d$.
We will focus on functions $h:\{0,1\}^{d'} \to \N$,
and prove the lower bound of $\Omega(d \log d)$ for this class of functions,
as $\Omega(d \log d) = \Omega(d' \log d')$.


We partition $\{0,1\}^\dd$ into $d$ subcubes based on the $\log_2 d$ most significant bits.
Specifically, for $i \in [d]$, the $\ord{i}$ subcube is defined as
\[C_i := \{x\in \{0,1\}^\dd \mid \dec(x_{d'}x_{d'-1}\cdots x_{d+1}) = i - 1\},\]
where $\dec(z) := \sum_{i = 1}^p z_i 2^{i-1}$ denotes the integer equivalent 
of the binary string $z_p z_{p-1} \ldots z_1$\,. %

Let $m:=d$. For ease of notation, we denote the set of subcube indices by $[m]$ and the set of dimensions in a subcube by $[d]$.
We use $i,j\in [m]$ to index subcubes,
and $a,b\in [d]$ to index dimensions.
We now define a collection of random variables, where each variable may depend on the previous ones:
\begin{itemize}
    \item The logarithm of the number of \emph{action dimensions}, $k$: a number picked uniformly at random from $\left[({1}/{2})\log_2 d \right]$.
    \item The set of all action dimensions, $R$: a uniformly random subset of $[d]$ of size $2^k$.
    \item The action dimension for each subcube, $r_i$: for each $i \in [m]$, $r_i$ is picked from $R$ uniformly and independently at random.
    \item The direction for each dimension (in $\Yes$ distribution), $\alpha_b$: for each $b \in [d]$, $\alpha_b$ is picked from $\{-1,+1\}$ uniformly and independently at random.
    (Technically, $\alpha_b$ will only be used for each $b \in R$. We define it for all $b \in [d]$ so that it is independent of $R$.)
    \item The direction of potential violations for each subcube (in $\No$ distribution), $\beta_i$: for each $i \in [m]$, $\beta_i$ is picked from $\{-1,+1\}$ uniformly and independently at random.
\end{itemize}


We use $\bT$ to refer to the entire collection
%
%
%
%
%
$(k, R, \{r_i \mid i \in [m]\}, \{\alpha_b \mid b \in [d]\}, \{\beta_i \mid i \in [m]\})$
%
of random variables.  %
We denote the tuple $(k,R,\{r_i \mid i \in [m]\})$ by $\bS$, also referred
%
%
to as the \emph{shared randomness} common to
the  %
%
%
%
distributions $\Yes$ and $\No$.
%
Given $\bT$, the distributions $\Yes$ and $\No$ generate
the  %
functions $f_{\bT}$
and $g_{\bT}$, respectively, %
where
\begin{align*}
f_{\bT}(x) &:=  \sum_{b \in [d'] \setminus R} x_b 3^b + \alpha_{r_i} x_{r_i}3^{r_i}, %
\\
g_{\bT}(x) &:=  \sum_{b \in [d'] \setminus R} x_b 3^b + \beta_i x_{r_i}3^{r_i} %
\end{align*}
and $i$ is the subcube containing $x$, \ie, $i = \dec(x_{d'}x_{d'-1}\cdots x_{d+1}) + 1$.


%
%
%
%
%
%
%
%



\subsection{The sign function and the distance to unateness for hard distributions} \label{sec:sgn-fn}
We need to analyze when functions in our hard distributions are increasing or decreasing in a specified dimension. More generally, since we are looking at comparison-based testers, we need to understand how $h(x)$ and $h(y)$ compare for any points $x,y$ in the hypercube $\{0,1\}^{\dd}$ for any function $h$ in the support of the hard distributions. To help us with that, we define the sign function $\sgn_h(x,y)$ and analyze its behavior on our hard distribution.

Recall that $\dec(z) := \sum_{i = 1}^p z_i 2^{i-1}$ denotes the integer equivalent of the binary string $z = z_p z_{p-1} \ldots z_1$.
We say that $x \order y$ if $\dec(x) < \dec(y)$. Note that $\order$ forms a total ordering
on $\{0,1\}^{\dd}$.

\begin{definition}
Given $x,y\in \{0,1\}^\dd$ such that $x \order y$, and a function $h:\{0,1\}^\dd\to \N$, define
\begin{align*}
\sgn_h(x,y) =
\begin{cases}
1 &\text{ if } h(x) < h(y), \\
0 &\text{ if } h(x) = h(y), \\
-1 &\text{ if } h(x) > h(y).
\end{cases}
\end{align*}
\end{definition}

%

Next, we show that for every function in the support of our hard distributions and two points $x \order y$ in different subcubes, $\sgn_h(x,y) = 1$.

For a distribution $\cD$, let $\supp(\cD)$ denote the support of $\cD$.

\begin{claim}\label{clm:inv}
For all $h \in \supp(\Yes) \cup \supp(\No), x \in C_i$ and $y \in C_j$ such that $i,j \in [m]$ and $i < j$, we have \[\sgn_h(x,y) = 1.\]
\end{claim}

\begin{proof}
Let $b$ be the most significant coordinate on which $x$ and $y$ differ, \ie, $x_b \neq y_b$. Since $x \in C_i, y \in C_j$ and $i < j$, we know that $x_b = 0, y_b = 1$ and $b > d$. For all $a > d$, the coefficient of $x_a$ in the definition of $h$ is $3^a$, irrespective of whether $h$ is $f_\bT$ or $g_\bT$.
Then,
\begin{align*}
h(x) &\leq \sum\limits_{a \in [\dd]} 3^a x_a \leq \sum\limits_{a > b} 3^a x_a + \sum\limits_{a < b} 3^a ;\\
h(y) &\geq \sum\limits_{a \geq b} 3^a y_a - \sum\limits_{a < b} 3^a.
\end{align*}
Thus, $h(y) - h(x) = 3^b - 2 \sum_{a < b} 3^a > 0$.
%
%
%
%
%
\end{proof}

%

Now we analyze the sign function on points $x,y$ in the same subcube. Its value is determined by the most significant coordinate on which $x$ and $y$ differ that has a nonzero coefficient in the definitions of hard functions $f_\bT$ and $g_\bT$. We define this coordinate next.

\begin{definition} \label{def:t}
For any setting of the shared randomness $\bS$,
index %
$i \in [m]$, and points $x,y \in C_i$, define $\coord{i}{\bS}(x,y)$
to be $\perp$ if $x$ and $y$ are the same on all coordinates in $([d] \setminus R) \cup \{r_i\}$; otherwise, define $\coord{i}{\bS}(x,y)$ to be the \emph{most significant coordinate}
in $([d] \setminus R) \cup \{r_i\}$ on which $x$ and $y$ differ.
\end{definition}

Note that $\bS$ determines $R$ and $\{r_i\}$.
For any $\bT$ that extends $\bS$, the restrictions of both $f_{\bT}$ and $g_{\bT}$ to $C_i$ is constant with respect to the coordinates in $R \setminus \{r_i\}$.
Thus, $\coord{i}{\bS}(x,y)$ is $\perp$ if $x$ and $y$ differ only on coordinates that do not influence the value of the function in $C_i$; otherwise, it is the first coordinate on which $x$ and $y$ differ that is influential
in $C_i$.

\begin{claim}\label{clm:interesting}
Fix some shared randomness $\bS$, 
index %
$i \in [m]$, and points $x,y \in C_i$ where $x \order y$.
Let $a = \coord{i}{\bS}(x,y)$.
For any $\bT$ that extends $\bS$,
\begin{itemize}
    \item if $a = \perp$, then $\sgn_{f_{\bT}}(x,y) = \sgn_{g_{\bT}}(x,y) = 0$;
    \item if $a \in ([d] \setminus R)$, then $\sgn_{f_{\bT}}(x,y) = \sgn_{g_{\bT}}(x,y) = 1$;
    \item if $a = r_i$, $\sgn_{f_{\bT}}(x,y) = \alpha_{a}$ and $\sgn_{g_{\bT}}(x,y) = \beta_i$.
\end{itemize}
\end{claim}

\begin{proof}
Recall that
%
\begin{align*}
f_{\bT}(x) &= \sum_{b \in [\dd] \setminus R} x_b 3^b + \alpha_{r_i} x_{r_i} 3^{r_i};\\
g_{\bT}(x) &= \sum_{b \in [\dd] \setminus R} x_b 3^b + \beta_i x_{r_i} 3^{r_i}.
\end{align*}

First, consider the case $a = \perp$. Then $x_b = y_b$ for all $b \in ([\dd] \setminus R) \cup \{r_i\}$. 
By the definition of $f_\bT$ and $g_\bT$\,, %
we get $f_\bT (x) = f_\bT (y)$ and $g_\bT (x) = g_\bT (y)$. Hence, $\sgn_{f_{\bT}}(x,y) = \sgn_{g_{\bT}}(x,y) = 0$.

Next, consider the case $a \neq r_i$\,. %
Then $a \notin R$ by the definition of $\coord{i}{\bS}(x,y)$.
Also, $x_b = y_b$ for all $b > a$ such that $b \notin R$.
Since $x \order y$, we have $x_a = 0$ and $y_a = 1$. Thus,
$f_{\bT}(y) - f_{\bT}(x) \geq 3^a - \sum_{b < a} 3^b > 0$ implying that $\sgn_{f_{\bT}}(x,y)=1$.
The same argument holds for $g_{\bT}$\,. %

Finally, consider the case $a = r_i$. Thus,
$f_{\bT}(y) - f_{\bT}(x) = \alpha_a 3^a + \sum_{b < a, b \notin R} (y_b - x_b) 3^b$.
Note that $-3^a < \sum_{b < a, b \notin R} (y_b - x_b) 3^b < 3^a$.
Hence, $f_{\bT}(y) - f_{\bT}(x)$ is positive when $\alpha_a = 1$ and negative when $\alpha_a = -1$. This implies that  $\sgn_{f_{\bT}}(x,y)=\alpha_a$\,. %
%
%
%
By an analogous argument,  $\sgn_{g_{\bT}}(x,y) = \beta_i$\,. %
\end{proof}

We complete \Sec{sgn-fn} by using Claims~\ref{clm:inv} and~\ref{clm:interesting} to analyze the distance to unateness for functions distributed according to our hard distributions.
\begin{claim}\label{clm:yes}
Every function $f_{\bT} \in\supp(\Yes)$ is unate.
\end{claim}

\begin{proof}
Fix some $f_{\bT} \in \supp(\Yes)$.
We need to prove that along each dimension $b \in [d']$, the function ${f_{\bT}}$ is either nonincreasing or nondecreasing.
Recall that a \emph{$b$-edge} is a pair $(x,y)$ of points in $\hypercube$, where $x_b=0,y_b=1$, and $x_a=y_a$ for all $a\in [\dd] \setminus \{b\}$.
Consider a $b$-edge $\{x,y\}$. If $b \in [d']\setminus R$, then, by Claims~\ref{clm:inv} and~\ref{clm:interesting}, $\sgn_{f_{\bT}}(x,y)=1,$ and hence ${f_{\bT}}$ is increasing in dimension~$b$.

If $b\in R$, then, by \Clm{interesting}, $\sgn_{f_{\bT}}(x,y)$ is 0 or $\alpha_b$\,. %
That is, if $\alpha_b=1$, then $f_{\bT}$ is nondecreasing, and if $\alpha_b=-1$, then $f_{\bT}$ is nonincreasing.
%
%
%
%
%
%
%
%
%
%
\end{proof}

We write $f \sim \cD$ to denote that $f$ is sampled from distribution $\cD$.

\begin{claim}\label{clm:no}
For sufficiently large $d$, a function $g_{\bT} \sim \No$ is $1/8$-far from unate with probability at least $19/20$.
\end{claim}

\begin{proof}
%
For each $r \in R$, let $A_r = \{ i \in [m] \mid r_i = r \}$ be the set of subcube indices with $r_i = r$.
Then
\begin{align}
\sum\limits_{r \in R} |A_r| = m = d \label{eqn:A_r}.
\end{align}
Hence, $\E[|A_r|] = d/|R| \geq \sqrt{d},$
since $|R| \leq \sqrt{d}$. By Chernoff and union bounds, for all $r \in R$,
we have $|A_r| \geq \sqrt{d}/2$	with probability at least $1 - d\exp(-\sqrt{d}/8)$.
	
Condition on the event that $|A_r| \geq \sqrt{d}/2$ for all $r \in R$.
Fix an $r \in R$.
For each $i \in A_r$\,, %
there is a random choice of $\beta_i$\,. %
Partition $A_r$ into $A^+_r$ and $A^-_r$,
depending on whether $\beta_i$ is
%
%
$+1$ or $-1$.   %
Again, by Chernoff and union bounds, for all $r \in R$, we have $\min(|A^+_r|,|A^-_r|) \geq |A_r|/4$ with probability at least $1-d\exp(-\sqrt{d}/32)$.

Thus, for sufficiently large $d$ and any choice of $R$,
\[
\Pr \left[\forall r \in R, \min(|A_r^+|, |A_r^-|) \geq \frac{|A_r|}{4} \right] \geq 1 - d\left(\exp \left(-\frac{\sqrt{d}}{8} \right) + \exp \left(-\frac{\sqrt{d}}{32} \right) \right) \geq \frac{19}{20}.
\]


%
Note that for all $i \in [m]$, the size $|C_i| = 2^d$.
By \Clm{interesting}, for any $r$-edge $(x,y)$ that lies within 
subcube $C_i$ where $i \in A_r$\,, %
the sign $\sgn_{g_{\bT}}(x,y) = \beta_i$\,. %
Hence, in $g_{\bT}$\,,  %
for all $i \in A^+_r$, all $r$-edges in subcube $C_i$ are increasing, whereas, for all $j \in A^-_r$, all $r$-edges in $C_j$ are decreasing.
In any unate function, all these $r$-edges are either nonincreasing or nondecreasing.
Thus $g_{\bT}$ differs from a unate function on at least
\[
  \frac{2^d}{2} \cdot \min(|A^+_r|,|A^-_r|) \geq \frac{2^d |A_r|}{8}
\]
points.  Overall, we need to change at least
$({2^d}/{8})\sum_{r \in R} |A_r|$ values. Since the $A_r$'s
partition the set of subcube indices,
$({2^d}/{8}) \sum_{r \in R} |A_r| = {2^{d'}}/{8}$.
\end{proof}

\subsection{Bad events}\label{sec:bad}
Fix a set of queries $Q$ made by a deterministic, nonadaptive, comparison-based tester.
We first identify certain \emph{bad} values of $\bS$,
%
on which $Q$ could potentially distinguish between $f_{\bT}$ and $g_{\bT}$ for any $\bT$ that extends $\bS$.
In this section, we prove that, for a given $Q$, the probability
of a bad $\bS$ is small.
%
In \Sec{indisting}, we show that the tester with queries $Q$ cannot distinguish between $f_{\bT}$ and $g_{\bT}$ for any $\bT$ that extends
%
good~$\bS$.
%

Recall that when a comparison-based tester queries points $x$ and $y$, it only sees $\sgn_h(x,y)$.  In order for $\sgn_h(x,y)$ to have a possibility of being different for $h\sim\Yes$ and $h\sim\No$, by Claims~\ref{clm:inv} and~\ref{clm:interesting}, queries $x$ and $y$ have to be in the same subcube $C_i$ for some $i\in [m]$ and, moreover, $\coord{i}{\bS}(x,y)$, the most significant coordinate in $([d] \setminus R) \cup \{r_i\}$ on which $x$ and $y$ differ has to be the action dimension $r_i$ for the subcube $C_i$. Subcubes are fixed, whereas action dimensions $R$ and $r_i$ are chosen randomly. Intuitively, the value of $\sgn_h(x,y)$ is likely to be determined by the top several coordinates in $[d]$ on which $x$ and $y$ differ. (And specifically, considering 5 top coordinates is sufficient to get high enough probability.) Next, we formalize this intuition in the definition of dimensions captured by $(x,y)$ and, more generally, by a set of query points.

\begin{definition}[Captured coordinates]\label{def:cap}
For points $x,y \in \hypercube$, define $\capt(x,y)$ to be
	\begin{enumerate}
	\item the set of all coordinates on which $x$ and $y$ differ if the size of the set is at most 5;
	\item the set of 5 most significant coordinates %
	on which $x$ and $y$ differ, otherwise.
	\end{enumerate}
We say that the pair $(x,y)$ \emph{captures} the set of coordinates $\capt(x,y)$.
For a set of points
$P\subseteq \{0,1\}^{\dd}$, define $\capt(P) := \bigcup_{x,y \in P} \capt(x,y)$
to be the set of all coordinates captured by the set $P$.
\end{definition}

For each $i \in [m]$, let $Q_i := Q \cap C_i$ denote the set of the queried points that lie in the subcube $C_i$\,. %
We define two bad events for $\bS$.
\begin{definition}[Bad events]\label{def:bad}
$\bS$ is \emph{bad} if at least one of the following bad events holds:
\begin{itemize}
\item Abort Event $\cA$:\ There exist $x,y \in Q$
such that $|\capt(x,y)| = 5$ and $\capt(x,y) \subseteq R$.
\item Collision Event $\cC$:\ There exist distinct $i,j \in [m]$
with $r_i = r_j$\,, %
such that $r_i\in \capt(Q_i)$ and \\ $r_j\in \capt(Q_j)$.
%
\end{itemize}
\end{definition}

If $\cA$ does not occur, then
for every pair $(x,y)$, the sign $\sgn_h(x,y)$ is determined by $\capt(x,y)$ for any $h\in \supp(\Yes)\cup \supp(\No)$.
And, if $\cC$ does not occur, conditioned on $\cA$ not occurring, then the tester cannot distinguish $\Yes$ and $\No.$
The heart of the analysis lies in proving that the bad events happen rarely.
%
%
%
%
%
%

\begin{lemma}\label{lem:abort}
Let $\lbconst = {1}/{20000}$. If $d$ is sufficiently large and $|Q|\leq \lbconst d\log d$, then $\Pr[\cA] < 0.01$\,. %
\end{lemma}

\begin{proof}
Recall that $k$ is the log of the number of action dimensions.
Fix any choice of $k$ (in $\bS$).
For each pair of points $x,y \in Q$ such that $|\capt(x,y)| = 5$,
\[\Pr[\capt(x,y) \subseteq R] = \frac{2^k}{d} \cdot \frac{2^k-1}{d-1} \cdots \frac{2^k-4}{d-4} < \Big(\frac{2^k}{d-4}\Big)^5.\]
Since $d-4 \geq d/2$ for all $d \geq 8$
and $k \leq (\log d)/2$, the probability is at most $32d^{-5/2}$.
By the union bound, $$\Pr[\cA] < |Q \times Q| \cdot 32d^{-5/2} < \lbconst^2 d^2 \log^2 d \cdot 32 d^{-5/2}  = 32 \lbconst^2 d^{-1/2} \log^2 d < 0.01$$ for all sufficiently large $d$.
\end{proof}

To analyze the probability of the collision event $\cC$, we prove the following combinatorial lemma that will be used to bound the number of coordinates $Q$ captures.

\begin{lemma}\label{lem:comb}
Let $c,d \in \mathbb{N}$ and $V$ be a set of $d$-dimensional vectors over an arbitrary alphabet. For all $x,y \in V$,
\begin{enumerate}
\item  if $x$ and $y$ differ on more than $c$ coordinates, then let $\capt_c(x,y)$ denote the set of the first $c$ coordinates on which $x$ and $y$ differ;
\item otherwise, let $\capt_c(x,y)$ denote the set of all coordinates
on which $x$ and $y$ differ.
\end{enumerate}
Define $\capt_c(V) := \bigcup_{x,y \in V} \capt_c(x,y)$.
Then $|\capt_c(V)| \leq c(|V|-1)$.
\end{lemma}

\begin{proof}
We construct $c$ different edge-colored graphs $G_1, \ldots, G_c$ over the vertex set $V$ with colors from $[d]$.
For every coordinate $i\in \capt_c(V)$, add one edge of color $i$ to exactly one of the graphs $G_t$\,. %
Since $i \in \capt_c(V)$, there exists at least one pair of
vectors $x,y$ such that $i \in \capt_c(x,y)$.
Thinking of each $\capt_c(x,y)$ as an ordered set, find one pair of vectors
$(x,y)$ where $i$ appears ``earliest'' in $\capt_c(x,y)$.
Let the position of $i$ in this $\capt_c(x,y)$ be denoted by $t$.
Add the edge $(x,y)$ to $G_t$ and color it $i$.
Note that each edge $(x,y)$ is added to $G_t$ at most once, and hence each graph $G_t$ is simple.

We claim that each $G_t$ is acyclic. Suppose not.
Let $C$ be a cycle in $G_t$\,. %
Let $(x,y)$ be the edge in $C$ with the smallest color $i$. 
Clearly, $x_i \neq y_i$\,, %
since $i \in \capt_c(x,y)$. There must exist another edge $(u,v)$ in $C$
such that $u_i \neq v_i$\,. %
Furthermore, the color of $(u,v)$
is $j > i$. Thus, $j$ is the $\ord{t}$ entry in $\capt_c(u,v)$.
Note that $i \in \capt_c(u,v)$,
since $u_i \neq v_i$ and $j > i$. It follows that $i$
appears earlier in $\capt_c(u,v)$ than in $\capt_c(x,y)$.
So the edge colored $i$ should not be in $G_t$\,, %
a contradiction.
\end{proof}

\Lem{comb} implies that a small $Q$ can capture only a few coordinates.
Next we bound the probability of the collision event.
%

\begin{lemma}\label{lem:collision}
If $|Q|\leq \lbconst d\log d$, then $\Pr[\cC] < 0.04$\,. %
%
\end{lemma}

\begin{proof}
By \Lem{comb},
\begin{align}
\sum_{i \in [m]} |\capt(Q_i)| \leq 5 \sum_{i \in [m]} |Q_i|  = 5|Q| \leq 5 \lbconst d\log d. \label{eqn:comb}
\end{align}

For each $r \in [d]$, define $A_r := \{j \in [m]\mid r\in \capt(Q_j)\}$ to be the set of indices of the subcubes in which coordinate $r$ is captured. Let $a_r := |A_r|$.
For each $\ell \in [\log d]$, define $n_\ell := |\{r \in [d]\mid a_r \in (2^{\ell-1},2^{\ell}]\}|$ to be the number of coordinates that are captured by more than $2^{\ell-1}$, but at most $2^\ell$ subcubes.
Observe that
\begin{align}
\sum_{\ell \in [\log d]} n_\ell 2^\ell \leq 2\sum_{r\in [d]} a_r = 2 \sum_{i \in [m]} |\capt(Q_i)| \leq 10 \lbconst d \log d, \label{eqn:006}
\end{align}
where the first inequality holds because if coordinate $r$ is included in the count $n_\ell$ then $2^{\ell-1}<a_r$, and the last inequality follows from \expeqref{Equation}{eqn:comb}.

Fix the log of the action dimensions, $k$, to be some $\kappa \in [({1}/{2})\log d ]$. For each $r\in [d]$, we say the event $\cC_r$ occurs if
 $r \in R$ and there exist distinct $i,j\in [m]$ such that $r_i = r_j = r$, and $r \in \capt(Q_i)\cap\capt(Q_j)$. By the union
bound, $\Pr[\cC \mid k = \kappa] \leq \sum_{r \in [d]} \Pr[\cC_r \mid k = \kappa]$.

Now, we compute $\Pr[\cC_r|k = \kappa]$.
Conditioning on $k = \kappa$, the size of the set of 
the %
action dimensions, $|R|$, is $2^{\kappa}$.
Only sets $Q_j$ with $j \in A_r$ capture coordinate $r$.
Event $\cC_r$ occurs if at least two of these sets have $r_i = r_j = r$. Hence,
\begin{align}
\Pr[\cC_r \mid k = \kappa] & = \Pr[r\in R]\cdot \Pr[\exists \text{ distinct } i,j\in A_r \text{ such that } r_i = r_j = r \mid r \in R] \label{eq:C_r-prob} \\
%
&= \frac{2^{\kappa}}{d} \cdot \sum\limits_{c = 2}^{a_r} \Pr\left[\text{There are exactly $c$ subcube indices } i\in A_r \text{ with } r_i=r \mid r \in R \right] \notag \\
& = \frac{2^{\kappa}}{d}\cdot \sum\limits_{c = 2}^{a_r} \binom{a_r}{c} \left(\frac{1}{2^{\kappa}} \right)^c \left(1-\frac{1}{2^{\kappa}} \right)^{a_r-c}, \label{eq:007}
\end{align}
where we used $\Pr[r \in R] = 2^\kappa/d$ to get the second equality, and $\Pr[r_i = r \mid r \in R] = 2^{-\kappa}$ for each $i \in [m]$ to get the third equality.
%
%

If $a_r > 2^{\kappa}/4$, by \expeqref{Equation}{eq:C_r-prob},
\begin{align}
\Pr[\cC_r \mid k = \kappa] \leq \Pr[r \in R] = 2^{\kappa}/d. \label{eq:big-a_r}
\end{align}

If $a_r \leq 2^{\kappa}/4$, by \expeqref{Equation}{eq:007},
\begin{align}
\Pr[\cC_r \mid k = \kappa] &= \frac{2^{\kappa}}{d}\cdot \sum\limits_{c = 2}^{a_r} \binom{a_r}{c} \left(\frac{1}{2^{\kappa}} \right)^c \left(1-\frac{1}{2^{\kappa}} \right)^{a_r-c} \notag \\
&< \frac{2^{\kappa}}{d} \left(1-\frac{1}{2^{\kappa}}\right)^{a_r} \sum_{c = 2}^{a_r} \left( a_r \cdot \frac{1}{2^{\kappa}} \cdot \left( 1-\frac{1}{2^{\kappa}} \right)^{-1} \right)^c  \notag \\
&\leq \frac{2^{\kappa}}{d} \sum_{c = 2}^{a_r} \left( \frac{a_r}{2^{\kappa-1}} \right)^c \notag \\
%
&< \frac{2^{\kappa}}{d} \cdot 2\left( \frac{a_r}{2^{\kappa-1}} \right)^2
= \frac{8a^2_r}{2^{\kappa} d}, \label{eq:small-a_r}
\end{align}
where we used
%
\[
  \binom{a_r}{c} < a_r^c \, , \quad (1 - 2^{-\kappa})^{-1} \leq
  2\quad\text{and}\quad \frac{a_r}{2^{\kappa-1}} \leq \frac{1}{2}
\]
%
%
to get the first, second and third inequalities, respectively.
%
Using
the   %
union bound over all $r \in [d]$ and grouping according to $n_\ell$,
we get
\begin{align}
\Pr[\cC \mid k = \kappa] \leq \sum_{r=1}^d \Pr[\cC_r \mid k = \kappa]
&= \sum_{r \in [d] : a_r > 2^\kappa/4} \Pr[\cC_r \mid k = \kappa] + \sum_{r \in [d]: a_r \leq 2^\kappa/4} \Pr[\cC_r \mid k = \kappa] \notag \\
&\leq \sum_{r \in [d] : a_r > 2^{\kappa-2}} \frac{2^{\kappa}}{d} + \sum_{r \in [d]: a_r \leq 2^{\kappa-2}} \frac{8a^2_r}{2^{\kappa} d} \notag \\
&\leq \frac{2^{\kappa}}{d} \sum\limits_{\ell = \max(\kappa - 1,1)}^{\log d} n_\ell + \frac{8}{d}  \sum_{\ell=1}^{\kappa-2} n_\ell 2^{2\ell - \kappa}, \label{eq:pr-c-k}
\end{align}
where the second last inequality follows from Equations~\eqref{eq:big-a_r} and~\eqref{eq:small-a_r}, and the last inequality holds because if coordinate $r$ is included in the count $n_\ell$, then $2^{\ell - 1} < a_r \leq 2^\ell$.

Averaging over all
the  %
values of $k$, we get
\begin{align}
\Pr[\cC] & = \frac{2}{\log d} \sum_{\kappa=1}^{\frac{\log d}{2}}\Pr[\cC | k = \kappa] \notag \\
& \leq \frac{16}{d \log d} \sum_{\kappa=1}^{\frac{\log d}{2}} \left(  2^{\kappa} \sum\limits_{\ell = \max(\kappa - 1,1)}^{\log d} n_\ell + \sum_{\ell=1}^{\kappa-2} n_\ell 2^{2\ell - \kappa} \right) \notag \\
%
%
%
& \leq \frac{16}{d\log d} \left(\sum_{\ell=1}^{\log d} n_\ell \sum\limits_{\kappa = 1}^{\ell+1} 2^{\kappa} + \sum_{\ell=1}^{\frac{\log d}{2} - 2} n_\ell \sum\limits_{\kappa = \ell + 2}^{\frac{\log d}{2}} 2^{2\ell - \kappa}  \right), \label{eq:pr-c}
\end{align}
where the first inequality follows from \expeqref{Equation}{eq:pr-c-k} and the last inequality is obtained by switching the order of summations and rearranging the terms.
Now,
\[
\sum_{\kappa = 1}^{\ell+1} 2^{\kappa} < 2^{\ell + 2}
\text{ and }
\sum_{\kappa= \ell+2}^{\frac{\log d}{2}} 2^{2\ell - \kappa} < 2^\ell\,.
\]
Using these inequalities to bound the sum in \expeqref{Equation}{eq:pr-c} and then applying \expeqref{Equation}{eqn:006}, we get
\[
\Pr[\cC] < \frac{80}{d\log d} \sum_{\ell=1}^{\log d} n_\ell 2^\ell
\leq \frac{80 \cdot 10 \lbconst d \log d}{d\log d}=0.04\,.\qedhere  %
\]
\end{proof}

\subsection{Indistinguishability of hard distributions}\label{sec:indisting}

\begin{proof}[{\bf Proof of \Thm{non-adap-lb}}]
Fix a deterministic, nonadaptive, comparison-based tester making a set of queries $Q$ whose size $|Q| \leq \lbconst d \log d$. (Recall that $\lbconst = {1}/{20000}$.)
The tester decides whether the input function $h$ is unate
based on the $\binom{|Q|}{2}$ comparisons between the function values at points in $Q$.
We define a labeled undirected graph $G^Q_h$ to be the clique on the node set $Q$, where each edge $\{x,y\}$ such that $x \order y$ is labeled by $\sgn_h(x,y)$.
The graph $G^Q_h$ defines the tester's ``view'' of the function $h$ after it has made its queries $Q$.

For any distribution $\cD$ over functions $f:\hypercube \to \N$, let $\cD$-view denote the distribution of the labeled graphs $G^Q_h$ when $h \sim \cD$.
For two distributions $\mathcal{D}_1$ and $\mathcal{D}_2$, let $\mathcal{D}_1 \approx_{\alpha} \mathcal{D}_2$ denote that the statistical distance between
$\mathcal{D}_1$ and $\mathcal{D}_2$ is at most $\alpha$.
By the version of Yao's principle stated in~\cite[Claim~5]{RS06},
it is sufficient to give two distributions, $\cD_1$ and $\cD_2$, 
on
%
%
positive and
%
%
negative instances, respectively, %
such that
the statistical distance between $\cD_1$-view and $\cD_2$-view is less than $1/6$
%
for every deterministic tester that makes at most $\lbconst d \log d$ queries.
Specifically, we show that for such testers, $\cD_1\text{-view} \approx_{3/19} \cD_2$-view.

We first show that, conditioned on the bad events $\cA$ and $\cC$ not occurring, the view of the tester is the same for both $\Yes$ and $\No$.
For a distribution $\cD$ and an event $E$, let $\cD|_{E}$ denote the distribution $\cD$ conditioned on the event $E$.

%
%
%
%
%

\begin{lemma}\label{lem:zero}
For every deterministic, nonadaptive, comparison-based tester,
\[\Yes\text{-view}|_{\overline{\cA \cup \cC}} = \No\text{-view}|_{\overline{\cA \cup \cC}}\,.\]  %
\end{lemma}
\begin{proof}
Fix some $\bS = S$ for $Q$ such that the bad events $\cA$ and $\cC$ do not occur. Let $\goodS$ denote the event that the collection of random variables $\bT$ used to define the support of $\Yes$ and $\No$ is extended from $S$.
We show that for any labeled clique $G$ over the vertex set $Q$,
\[\Pr_{f\sim \Yes|_{\goodS}} [G^Q_f = G] = \Pr_{g\sim \No|_{\goodS}} [G^Q_g = G].\]

Let $F = \supp(\Yes|_{\goodS}) \cup \supp(\No|_{\goodS})$.
Fix some $G \in \{ G^Q_h \mid h \in F \}$. Consider an edge $\{x,y\} \in G$ with $x \order y$ and $x,y \in \hypercube$.
If $x$ and $y$ lie in different subcubes, then, by \Clm{inv}, its label $\sgn_h(x, y) = 1$ for all $h \in F$. Similarly, if $x$ and $y$ lie in the same subcube $C_i$ for some $i \in [m]$, 
but with $\coord{i}{\bS}(x,y) \neq r_i$\,, %
then by \Clm{interesting}, $\sgn_h(x,y)$ is the same for all $h \in F$.
Hence the label of the edge $\{x,y\}$ can potentially distinguish between $\Yes$ and $\No$ only if $x$ and $y$ lie in the same subcube $C_i$ for some $i \in [m]$ and $\coord{i}{\bS}(x,y) = r_i$\,. %
We call such edges \emph{interesting}. In the rest of this proof, we focus on interesting edges, since the labels of non-interesting edges cannot be used to distinguish between $\Yes$ and $\No$.

Fix $i \in [m]$. Recall that $Q_i = Q \cap C_i$  is the set of queried points that lie in the subcube $C_i$\,. %
Let $G(Q_i)$ denote the subgraph of $G$ induced on vertex set $Q_i$\,. %
By \Clm{interesting}, all interesting edges in $G(Q_i)$ have the same label. 
Denote this label by $\ell_i$\,. %

Let $I = \{ i \in [m] \mid G(Q_i) \text{ contains an interesting edge}\}$ denote the set of indices of all subcubes with at least one interesting edge.

Now, we focus on $g \sim \No|_{\goodS}$\,. %
For any edge $\{x,y\} \in G$ with $x \order y$, let $\ell(x,y)$ denote its label. The probability
\begin{align}
\Pr\limits_{g \sim \No|_{\goodS}}[G^Q_g = G] &= \Pr \Big[ \bigwedge_{i \in I} \bigwedge_{\substack{\{x,y\} \in G(Q_i), \\ \{x, y \} \text{ is interesting}}} \ell(x,y) = \sgn_g (x,y) \Big] \notag \\
&= \Pr \Big[ \bigwedge_{i \in I} \bigwedge_{\substack{\{x,y\} \in G(Q_i), \\ \{x, y \} \text{ is interesting}}} \ell(x,y) = \beta_i \Big] \tag*{(by \Clm{interesting})} \\
&= \Pr \Big[ \bigwedge_{i \in I}  (\ell_i = \beta_i) \Big] \notag \\
&= \prod\limits_{i \in I}\Pr[\ell_i = \beta_i] = 2^{-|I|}, \label{eq:no-view-graph}
\end{align}
since each $\beta_i$ is chosen uniformly and independently at random from $\{+1, -1\}$.
%

%
%
Similarly,
for $f \sim \Yes|_{\goodS}$\,, 
%
the probability
\begin{align}
\Pr\limits_{f \sim \Yes|_{\goodS}}[G^Q_f = G] = \Pr \Big[ \bigwedge_{i \in I}  (\ell_i = \alpha_{r_i}) \Big]. \label{eq:yes-view-graph}
\end{align}
Note that unless there exist distinct $i,j \in I$ 
such that $r_i = r_j$\,, %
the individual events in the probability expression \expeqref{Equation}{eq:yes-view-graph} are mutually independent.
We show that if the bad events $\cA$ and $\cC$ do not occur then all subcubes $C_i$ with interesting edges 
have distinct action dimensions $r_i$\,. %

\begin{claim}\label{clm:good-events}
If $\overline{\cA \cup \cC}$ then $r_i \neq r_j$ for all distinct $i,j \in I$.
\end{claim}
\begin{proof}
Assume for the sake of contradiction that there exist distinct $i,j \in I$ 
with $r_i = r_j$\,. %
The fact that $\overline{\cC}$ holds ensures that if $r_i \in \capt(Q_i)$ and $r_j \in \capt(Q_j)$ for distinct $i, j \in [m]$, 
then $r_i \neq r_j$\,. %
Hence, $r_i \notin \capt(Q_i)$ or $r_j \notin \capt(Q_j)$. Suppose without loss of generality $r_i \notin \capt(Q_i)$.
Since $i \in I$, there must exist an interesting edge $\{x,y\} \in G(Q_i)$ with $x \order y$. We know $r_i \notin \capt(Q_i)$, so $r_i \notin \capt(x,y)$. By \Def{cap}, if $x$ and $y$ differ on at most five coordinates, then $r_i \in \capt(x,y)$. Hence, $x$ and $y$ differ on at least six coordinates, and $r_i$ is not among the five most significant coordinates. Since the abort event $\cA$ does not occur, $\capt(x,y) \nsubseteq R$. Hence, there exists a coordinate $a \in \capt(x,y)$ that is not in $R$ and is 
more significant than $r_i$\,. %
Then, $\coord{i}{\bS}(x,y) \geq a$, contradicting the fact that $\coord{i}{\bS}(x,y) = r_i$\,. %
Hence, $r_i \neq r_j$ for all distinct $i,j \in I$.
\end{proof}

By \expeqref{Equation}{eq:yes-view-graph} and \Clm{good-events},
\begin{align}
\Pr\limits_{f \sim \Yes|_{\goodS}}[G^Q_f = G] = \prod\limits_{i \in I} \Pr[\ell_i = \alpha_{r_i}] = 2^{-|I|}, \label{eq:yes-bound}
\end{align}
since $\alpha_b$ is chosen uniformly and independently at random from $\{-1, +1\}$ for each $b \in [d]$. By Equations~\eqref{eq:no-view-graph} and~\eqref{eq:yes-bound},
\[\Pr_{f\sim \Yes|_{\goodS}} [G^Q_f = G] = \Pr_{g\sim \No|_{\goodS}} [G^Q_g = G],\]
%
%
completing the proof of \Lem{zero}.
%
\end{proof}

We now wrap up the proof of \Thm{non-adap-lb}.
Let $\Far$ denote the event that a function $h \sim \No$ is ${1}/{8}$-far from unate.
Define $\FarNo$ to be $\No|_\Far$\,. %
Note that every function $h \sim \FarNo$ is ${1}/{8}$-far from unate. By \Clm{yes}, every function $h \sim \Yes$ is unate. Hence, to complete the proof of this theorem, it suffices to show that $\Yes\text{-view} \approx_{3/19} \FarNo$-view.

By \Clm{no}, $\Pr[\Far] \geq 19/20$.
We use the following claim from~\cite{RS06} %
to show that $\No$ and $\FarNo$ are close.

\begin{claim}[Claim 4 in \cite{RS06}]\label{clm:sd-prob}
Let $E$ be an event that happens with probability at least $\alpha = 1 - 1/a$ under the distribution $\cD$.
Then, $\cD|_{E} \approx_{\alpha'} \cD$ where $\alpha' = 1/(a-1)$.
\end{claim}

%
By \Clm{sd-prob},    %
$\No \approx_{1/19} \FarNo$. Consequently,
\begin{align}
\No\text{-view} \approx_{1/19} \FarNo\text{-view}. \label{eq:no-farno}
\end{align}

By Lemmas~\ref{lem:abort} and~\ref{lem:collision}, the probability $\Pr[\overline{\cA \cup \cC}] \geq 19/20$. Using
%
\Clm{sd-prob}, we get   %
%
\begin{align}
\Yes\text{-view} &\approx_{1/19} \Yes\text{-view}|_{\overline{\cA \cup \cC}}\,, \label{eq:yes-view}\\
\No\text{-view} &\approx_{1/19} \No\text{-view}|_{\overline{\cA \cup \cC}}\,. \label{eq:no-view}
\end{align}
%

Using  \expeqref{Equation}{eq:yes-view}, \Lem{zero}, \expeqref{Equation}{eq:no-view}, and, finally, \expeqref{Equation}{eq:no-farno}, we get
\[\Yes\text{-view} \approx_{1/19} \Yes\text{-view}|_{\overline{\cA \cup \cC}} =  \No\text{-view}|_{\overline{\cA \cup \cC}} \approx_{1/19} \No\text{-view} \approx_{1/19} \FarNo\text{-view}.\]
Hence $\Yes\text{-view} \approx_{3/19} \FarNo\text{-view}$, completing the proof of \Thm{non-adap-lb}.
\end{proof}


\section{Other lower bounds}
\subsection{The lower bound for adaptive testers over the hypergrid} \label{sec:adap-lb}
We show that every unateness tester for functions $f:[n]^d \to \R$  requires
\[
  \Omega\left(\frac{d \log n}{\eps}-\frac{\log 1/\eps}{\eps}\right)
\]
queries for $\eps\in(0,1/4)$ and prove \Thm{adap-lb}.

\begin{proof}[Proof of \Thm{adap-lb}]
  By Yao's minimax principle and the reduction to testing with comparison-based testers from~\cite{Fis04,CS14} (stated for completeness in \Thm{CS14}), it is sufficient to give a hard input distribution on which every deterministic comparison-based tester fails with probability at least $2/3$. We use the hard distribution constructed  by Chakrabarty and Seshadhri~\cite{CS14}, who used it to prove the same lower bound for testing monotonicity. Their distribution is a mixture of two distributions, $\Yes$ and $\No,$ on positive and negative instances, respectively.
%
%
%
The positive %
instances for their problem are
%
%
%
%
functions that are monotone and, therefore, unate;
the %
negative instances are
%
%
functions that are $\eps$-far from monotone. We show that their 
%
%
distribution $\No$
%
is supported on functions that are $\eps$-far from unate, \ie, negative instances for our problem. Then the required lower bound for unateness follows from the fact that every deterministic comparison-based tester needs the stated number of queries to distinguish
the  %
%
%
distributions $\Yes$ and $\No$
%
with high enough probability.

We start by describing 
%
%
the distributions $\Yes$ and $\No$
%
used in \cite{CS14}. We will define them as distributions on functions over the hypercube domain. Next, we explain how to convert functions over hypercubes  to functions over hypergrids.

Without loss of generality, assume $n$ is a power of $2$ and let $\ell := \log_2 n$.
%
For all $z \in [n]$, let $\bin(z)$ denote the binary representation of $z-1$ as an $\ell$-bit vector $(z_1, \ldots, z_\ell)$, where $z_1$ is the least significant bit.

We now describe the mapping used to convert functions on hypergrids to functions on hypercubes. Let $\phi:[n]^d \to \{0,1\}^{d\ell}$ be the mapping that takes $y\in[n]^d$ to the concatenation of $\bin(y_1),\dots,\bin(y_d)$. Any function $f:\{0,1\}^{d \ell} \to \R$ can be easily converted into a function $\widetilde{f}:[n]^d \to \R$, where $\widetilde{f}(y) := f(\phi(y))$.

Let $m:= d\ell$. For $x \in \{0,1\}^m$, recall that $\dec(x) = \sum\nolimits_{i=1}^m x_i 2^{i-1}$ denotes the integer equivalent of the binary number represented by vector $x$. For simplicity, assume $1/\eps$ is a power of $2$.
Partition the set of points $x \in \{0,1\}^m$ according to the most significant $\log(1/2\eps)$ dimensions.
That is, for $k \in [1/2\eps]$, let
\[S_k := \{x: \dec(x) \in [(k-1)\cdot \eps 2^{m+1}, k\cdot\eps 2^{m+1} - 1]\}.\]
The hypercube is partitioned into $1/2\eps$ sets $S_k$ of equal size, and each $S_k$ forms a subcube of dimension $m' = m - \log(1/\eps) + 1$.

We now describe the distributions $\Yes$ and $\No$ for functions on hypercubes.
The distribution $\Yes$ consists of a single function $f(x) = 2 \dec(x)$.
The distribution $\No$ is the uniform distribution over $m'/2\eps$ functions $g_{j,k}$, where $j \in [m']$ and $k \in [1/2\eps]$, defined as follows:
\begin{align*}
    g_{j,k}(x) =
    \begin{cases}
        2\dec(x) - 2^j - 1 &\text{ if } x_j = 1 \text{ and } x \in S_k;\\
        2\dec(x)
        %
        &\text{ otherwise.}
    \end{cases}
\end{align*}
To get the distributions $\Yes$ and $\No$ for the hypergrid, we convert $f$ to $\widetilde{f}$ and each function $g_{j,k}$ to $\widetilde{g_{j,k}}$\,, %
using the transformation defined before.


Chakrabarty and Seshadhri~\cite{CS14} proved that $f$ is monotone and each function $\widetilde{g_{j,k}}$  is $\eps$-far from monotone. It remains to show that functions $\widetilde{g_{j,k}}$ are also $\eps$-far from unate.

\begin{claim}
Each function $\widetilde{g_{j,k}}$ is $\eps$-far from unate.
\end{claim}

\begin{proof}
To prove that $\widetilde{g_{j,k}}$ is $\eps$-far from unate, it suffices to show that there exists a dimension $i$, such that there are at least $\eps 2^{d\ell}$ increasing $i$-pairs and at least $\eps 2^{d\ell}$ decreasing $i$-pairs w.r.t.\ $\widetilde{g_{j,k}}$ and that all of these $i$-pairs are disjoint. Let $u,v \in [n]^d$ be two points such that $\phi(u)$ and $\phi(v)$ differ only in the $\ord{j}$ bit. Clearly, $u$ and $v$ form an $i$-pair, where $i =\ceil{j/\ell}$. Now, if $\phi(u),\phi(v) \in S_k$ and $u \prec v $, then $\widetilde{g_{j,k}}(v) = \widetilde{g_{j,k}}(u) - 1$. So, the $i$-pair $(u,v)$ is decreasing. The total number of such $i$-pairs is  $2^{d\ell - \log(1/2\eps) - 1} = \eps 2^{d\ell}$. If $\phi(u),\phi(v) \in S_{k'}$ where $k' \neq k$, then the $i$-pair $(u,v)$ is increasing. Clearly, there are at least $\eps 2^{d\ell}$ such $i$-pairs. All the $i$-pairs we mentioned are disjoint.
Hence, $\widetilde{g_{j,k}}$ is $\eps$-far from unate.
\end{proof}
This completes the proof of \Thm{adap-lb}.
\end{proof}

\subsection{The lower bound for nonadaptive testers over the hypergrid}\label{sec:na-lb-hg}
The lower bound for nonadaptive testers over hypergrids follows from a combination of the lower bound for nonadaptive testers over hypercube and the lower bound for adaptive testers over hypergrids.

\begin{proof}[Proof of~\Thm{na-lb-hg}]
Fix $\eps = 1/8$. The proof consists of two parts. The lower bound for adaptive testers is also a lower bound for nonadaptive 
testers, %
and so, the bound of $\Omega(d \log n)$ holds.
Next, we extend the $\Omega(d \log d)$ lower bound for hypercubes.
Assume $n$ to be even. Define function $\psi:[n] \to \{0,1\}$ as $\psi(a):= \mathbbm{1}[a > n/2]$ for $a \in [n]$. For $x = (x_1, x_2, \ldots, x_d) \in [n]^d$, define the mapping $\Psi: [n]^d \to \{0,1\}^d$ as
\[\Psi(x) := (\psi(x_1), \psi(x_2), \ldots, \psi(x_d)).\]
Any function $f:\{0,1\}^d \to \R$ can be extended to $\tilde{f}:[n]^d \to \R$ using the mapping $\tilde{f}(x) = f(\Psi(x))$ for all $x \in [n]^d$. The proof of \Thm{non-adap-lb} goes through for hypergrids as well, and so we have an $\Omega(d \log d)$ lower bound. Combining the two lower bounds, we get a bound of $\Omega(d \cdot \max\{\log n, \log d\})$, which is asymptotically equal to $\Omega(d(\log n + \log d))$.
\end{proof}

\section{Conclusion and open questions}
In this work, we give the first algorithms for
testing unateness of real-valued functions over the hypercube as well as the hypergrid domains.
We also show that our algorithms are optimal by proving matching lower bounds, thus resolving the query complexity of testing unateness of real-valued functions.
Our results demonstrate that, for real-valued functions, adaptivity helps with testing unateness, which is not the case in monotonicity testing.

%

Subsequent to the initial publication of this work~\cite{BCPRS17a}, the problem of testing unateness of Boolean functions received significant attention.
Concurrent with our work, Chen et al.~\cite{CWX17} proved a lower bound of $\Omega\left({d^{2/3}}/{\log^3 d} \right)$ for adaptive unateness testers of Boolean functions over $\{0,1\}^d$.
Subsequently, Chen et al.~\cite{CWX17focs} gave an adaptive unateness tester with query complexity $\widetilde{O}\left({d^{3/4}}/{\eps^2}\right)$ for the same class of functions.
An exciting recent development is an $\widetilde{O}\left({d^{2/3}}/{\eps^2}\right)$-query algorithm for this problem by Chen and Waingarten~\cite{CW19}, which only leaves a polylogarithmic (in $d$) gap between the upper bound and the lower bound.
%

Next, we discuss nonadaptive unateness testing of Boolean functions over $\{0,1\}^d$. In a subsequent work, Baleshzar et al.~\cite{BCPRS17b} proved a lower bound of $\Omega\left({d}/{\log d} \right)$ for one-sided error testers for this problem.
Since Boolean functions are a special case of real-valued functions,
our nonadaptive algorithm over the hypercube also works for Boolean functions.
This algorithm has 1-sided error and its query complexity  is $O\left(({d}/{\eps})\log ({d}/{\eps})\right)$ which is currently the best known upper bound for any (even 2-sided error) nonadaptive algorithm. There is still a polylogarithmic (in $d$) gap between the upper bound and the lower bound.
An interesting open question is to determine if testers with two-sided error have better query complexity than testers with one-sided error in the nonadaptive setting.
%

Finally, we mention some recent work on approximating the distance to unateness of Boolean functions over the hypercube domain. Levi and Waingarten~\cite{LW19} showed that every algorithm approximating the distance to unateness within a constant factor requires $\widetilde{\Omega}(d)$ queries and strengthened their lower bound to $\widetilde{\Omega}(d^{3/2})$ queries for nonadaptive algorithms. Subsequently, Pallavoor et al.~\cite{PRW20} proved that every nonadaptive algorithm approximating the distance to unateness within a $d^{1/2 - k}$ factor for $k > 0$ requires $2^{d^k}$ queries. No nontrivial upper bounds are currently known for this problem.

\paragraph{Acknowledgments.} We thank Oded Goldreich for useful discussions and Meiram Murzabulatov for participation in initial discussions on this work.
%
We would like to thank \emph{Theory of Computing} and L\'aszl\'o Babai 
personally
for the unparalleled copy editing service they provided. We have never seen
anything like this from other journals. The comments and suggestions we
received helped us greatly improve our paper. We are also pleased with its
appearance. 


%% Theory of Computing
%% File: v016a003.bbl
%% Publication date: September 16, 2020
%% Authors: Roksana Baleshzar, Deeparnab Chakrabarty, Ramesh Krishnan S. Pallavoor, Sofya Raskhodnikova, and C. Seshadhri
\newcommand{\prelim}{Preliminary version in\ }
% Define empty bibhead if not already defined
\providecommand{\bibhead}[1]{}
% Define tocrefpdfbookmark if not already defined
\expandafter\ifx\csname pdfbookmark\endcsname\relax%
  \providecommand{\tocrefpdfbookmark}{}
\else\providecommand{\tocrefpdfbookmark}{%
   \hypertarget{tocreferences}{}%
   \pdfbookmark[1]{References}{tocreferences}}%
\fi

\tocrefpdfbookmark
\begin{thebibliography}{10}

\bibitem{AC06}\bibhead{AC06}
{\sc Nir Ailon and Bernard Chazelle}: Information theory in property testing
  and monotonicity testing in higher dimension.
\newblock {\em Inform. and Comput.}, 204(11):1704--1717, 2006.
\newblock \prelim
  \href{https://doi.org/10.1007/978-3-540-31856-9\_36}{STACS'05}.
\newblock [\epfmtdoi{10.1016/j.ic.2006.06.001}]

\bibitem{AJMR16}\bibhead{AJMR16}
{\sc Pranjal Awasthi, Madhav Jha, Marco Molinaro, and Sofya Raskhodnikova}:
  Testing {Lipschitz} functions on hypergrid domains.
\newblock {\em Algorithmica}, 74(3):1055--1081, 2016.
\newblock \prelim
  \href{http://doi.org/10.1007/978-3-642-32512-0\_33}{RANDOM'12}.
\newblock [\epfmtdoi{10.1007/s00453-015-9984-y}]

\bibitem{BCPRS17b}\bibhead{BCPRS17b}
{\sc Roksana Baleshzar, Deeparnab Chakrabarty, Ramesh Krishnan~S. Pallavoor,
  Sofya Raskhodnikova, and C.~Seshadhri}: A lower bound for nonadaptive,
  one-sided error testing of unateness of {Boolean} functions over the
  hypercube.
\newblock {\em Electron. Colloq. on Comput. Complexity (ECCC)}, 2017.
\newblock [\epfmt{eccc}{TR17-111}, \epfmt{arxiv}{1706.00053}]

\bibitem{BCPRS17a}\bibhead{BCPRS17a}
{\sc Roksana Baleshzar, Deeparnab Chakrabarty, Ramesh Krishnan~S. Pallavoor,
  Sofya Raskhodnikova, and C.~Seshadhri}: Optimal unateness testers for
  real-valued functions: Adaptivity helps.
\newblock In {\em Proc. 44th Internat. Colloq. on Automata, Languages and
  Programming (ICALP'17)}, pp. 5:1--5:14. Springer, 2017.
\newblock [\epfmtdoi{10.4230/LIPIcs.ICALP.2017.5}, \epfmt{arxiv}{1703.05199}]

\bibitem{BRW05}\bibhead{BRW05}
{\sc Tu\u{g}kan Batu, Ronitt Rubinfeld, and Patrick White}: Fast approximate
  {PCP}s for multidimensional bin-packing problems.
\newblock {\em Inform. and Comput.}, 196(1):42--56, 2005.
\newblock \prelim
  \href{https://doi.org/10.1007/978-3-540-48413-4\_25}{RANDOM'99}.
\newblock [\epfmtdoi{10.1016/j.ic.2004.10.001}]

\bibitem{BCHKS96}\bibhead{BCHKS96}
{\sc Mihir Bellare, Don Coppersmith, Johan H{\aa}stad, Marcos~A. Kiwi, and
  Madhu Sudan}: Linearity testing in characteristic two.
\newblock {\em IEEE Trans. Inform. Theory}, 42(6):1781--1795, 1996.
\newblock \prelim \href{http://doi.org/10.1109/SFCS.1995.492574}{FOCS'95}.
\newblock [\epfmtdoi{10.1109/18.556674}]

\bibitem{B18}\bibhead{B18}
{\sc Aleksandrs Belovs}: Adaptive lower bound for testing monotonicity on the
  line.
\newblock In {\em Proc. 22nd Internat. Workshop on Randomization and
  Computation (RANDOM'18)}, pp. 31:1--31:10. Schloss Dagstuhl--Leibniz-Zentrum
  fuer Informatik, 2018.
\newblock [\epfmtdoi{10.4230/LIPIcs.APPROX-RANDOM.2018.31},
  \epfmt{arxiv}{1801.08709}]

\bibitem{BB15}\bibhead{BB15}
{\sc Aleksandrs Belovs and Eric Blais}: Quantum algorithm for monotonicity
  testing on the hypercube.
\newblock {\em Theory of Computing}, 11(16):403--412, 2015.
\newblock [\epfmtdoi{10.4086/toc.2015.v011a016}, \epfmt{arxiv}{1503.02868}]

\bibitem{BB16}\bibhead{BB16}
{\sc Aleksandrs Belovs and Eric Blais}: A polynomial lower bound for testing
  monotonicity.
\newblock In {\em Proc. 48th STOC}, pp. 1021--1032. ACM Press, 2016.
\newblock [\epfmtdoi{10.1145/2897518.2897567}, \epfmt{arxiv}{1511.05053}]

\bibitem{BCLR08}\bibhead{BCLR08}
{\sc Michael Ben{-}Or, Don Coppersmith, Michael Luby, and Ronitt Rubinfeld}:
  Non-abelian homomorphism testing, and distributions close to their
  self-convolutions.
\newblock {\em Random Structures Algorithms}, 32(1):49--70, 2008.
\newblock \prelim
  \href{https://doi.org/10.1007/978-3-540-27821-4\_25}{RANDOM'04}.
\newblock [\epfmtdoi{10.1002/rsa.20182}]

\bibitem{BerRY14}\bibhead{BerRY14}
{\sc Piotr Berman, Sofya Raskhodnikova, and Grigory Yaroslavtsev}:
  {$L_p$-testing}.
\newblock In {\em Proc. 46th STOC}, pp. 164--173. ACM Press, 2014.
\newblock [\epfmtdoi{10.1145/2591796.2591887}]

\bibitem{BGJRW12}\bibhead{BGJRW12}
{\sc Arnab Bhattacharyya, Elena Grigorescu, Kyomin Jung, Sofya Raskhodnikova,
  and David~P. Woodruff}: Transitive-closure spanners.
\newblock {\em SIAM J. Comput.}, 41(6):1380--1425, 2012.
\newblock \prelim \href{https://doi.org/10.1137/1.9781611973068.101}{SODA'09}.
\newblock [\epfmtdoi{10.1137/110826655}, \epfmt{arxiv}{0808.1787}]

\bibitem{BCS18}\bibhead{BCS18}
{\sc Hadley Black, Deeparnab Chakrabarty, and C.~Seshadhri}: A $o(d)\cdot$
  polylog$(n)$ monotonicity tester for {Boolean} functions over the hypergrid
  $[n]^d$.
\newblock In {\em Proc. 29th ACM-SIAM Symp. on Discrete Algorithms (SODA'18)},
  pp. 2133--2151. ACM Press, 2018.
\newblock [\epfmtdoi{10.1137/1.9781611975031.139}, \epfmt{arxiv}{1710.10545}]

\bibitem{BCS20}\bibhead{BCS20}
{\sc Hadley Black, Deeparnab Chakrabarty, and C.~Seshadhri}: Domain reduction
  for monotonicity testing: {A} $o(d)$ tester for {Boolean} functions in
  $d$-dimensions.
\newblock In {\em Proc. 31st ACM-SIAM Symp. on Discrete Algorithms (SODA'20)},
  pp. 1975--1994. ACM Press, 2020.
\newblock [\epfmtdoi{10.1137/1.9781611975994.122}, \epfmt{arxiv}{1811.01427}]

\bibitem{BB17}\bibhead{BB17}
{\sc Eric Blais and Abhinav Bommireddi}: Testing submodularity and other
  properties of valuation functions.
\newblock In {\em Proc. 8th Innovations in Theoret. Computer Science
  (ITCS'17)}, pp. 33:1--33:17. Schloss Dagstuhl--Leibniz-Zentrum fuer
  Informatik, 2017.
\newblock [\epfmtdoi{10.4230/LIPIcs.ITCS.2017.33}, \epfmt{arxiv}{1611.07879}]

\bibitem{BBM12}\bibhead{BBM12}
{\sc Eric Blais, Joshua Brody, and Kevin Matulef}: Property testing lower
  bounds via communication complexity.
\newblock {\em Comput. Complexity}, 21(2):311--358, 2012.
\newblock \prelim \href{http://doi.org/10.1109/CCC.2011.31}{CCC'11}.
\newblock [\epfmtdoi{10.1007/s00037-012-0040-x}]

\bibitem{BlaRY14}\bibhead{BlaRY14}
{\sc Eric Blais, Sofya Raskhodnikova, and Grigory Yaroslavtsev}: Lower bounds
  for testing properties of functions over hypergrid domains.
\newblock In {\em Proc. 29th IEEE Conf. on Computational Complexity (CCC'14)},
  pp. 309--320. IEEE Comp. Soc. Press, 2014.
\newblock [\epfmtdoi{10.1109/CCC.2014.38}]

\bibitem{BLR93}\bibhead{BLR93}
{\sc Manuel Blum, Michael Luby, and Ronitt Rubinfeld}: Self-testing/correcting
  with applications to numerical problems.
\newblock {\em J. Comput. System Sci.}, 47(3):549--595, 1993.
\newblock \prelim \href{https://doi.org/10.1145/100216.100225}{STOC'90}.
\newblock [\epfmtdoi{10.1016/0022-0000(93)90044-W}]

\bibitem{BCGM12}\bibhead{BCGM12}
{\sc Jop Bri{\"{e}}t, Sourav Chakraborty, David Garc{\'{\i}}a{-}Soriano, and
  Arie Matsliah}: Monotonicity testing and shortest-path routing on the cube.
\newblock {\em Combinatorica}, 32(1):35--53, 2012.
\newblock \prelim
  \href{https://doi.org/10.1007/978-3-642-15369-3\_35}{RANDOM'10}.
\newblock [\epfmtdoi{10.1007/s00493-012-2765-1}]

\bibitem{C16a}\bibhead{C16a}
{\sc Deeparnab Chakrabarty}: Monotonicity testing.
\newblock In {\em Encyclopedia of Algorithms}, pp. 1352--1356. Springer, 2016.
\newblock [\epfmtdoi{10.1007/978-1-4939-2864-4\_699}]

\bibitem{CDJS17}\bibhead{CDJS17}
{\sc Deeparnab Chakrabarty, Kashyap Dixit, Madhav Jha, and C.~Seshadhri}:
  Property testing on product distributions: Optimal testers for bounded
  derivative properties.
\newblock {\em ACM Trans. Algorithms}, 13(2):20:1--20:30, 2017.
\newblock \prelim \href{https://doi.org/10.1137/1.9781611973730.121}{SODA'15}.
\newblock [\epfmtdoi{10.1145/3039241}, \epfmt{arxiv}{1404.0718}]

\bibitem{CS13}\bibhead{CS13}
{\sc Deeparnab Chakrabarty and C.~Seshadhri}: Optimal bounds for monotonicity
  and {Lipschitz} testing over hypercubes and hypergrids.
\newblock In {\em Proc. 45th STOC}, pp. 419--428. ACM Press, 2013.
\newblock [\epfmtdoi{10.1145/2488608.2488661}, \epfmt{arxiv}{1204.0849}]

\bibitem{CS14}\bibhead{CS14}
{\sc Deeparnab Chakrabarty and C.~Seshadhri}: An optimal lower bound for
  monotonicity testing over hypergrids.
\newblock {\em Theory of Computing}, 10(17):453--464, 2014.
\newblock \prelim
  \href{https://doi.org/10.1007/978-3-642-40328-6\_30}{RANDOM'13}.
\newblock [\epfmtdoi{10.4086/toc.2014.v010a017}]

\bibitem{CS16}\bibhead{CS16}
{\sc Deeparnab Chakrabarty and C.~Seshadhri}: An $o(n)$ monotonicity tester for
  {Boolean} functions over the hypercube.
\newblock {\em SIAM J.\ Comput.}, 45(2):461--472, 2016.
\newblock \prelim \href{https://doi.org/10.1145/2488608.2488660}{STOC'13}.
\newblock [\epfmtdoi{10.1137/13092770X}, \epfmt{arxiv}{1302.4536}]

\bibitem{CS19}\bibhead{CS19}
{\sc Deeparnab Chakrabarty and C.~Seshadhri}: Adaptive {Boolean} monotonicity
  testing in total influence time.
\newblock In {\em Proc. 10th Innovations in Theoret. Computer Science
  (ITCS'19)}, pp. 20:1--20:7. Schloss Dagstuhl--Leibniz-Zentrum fuer
  Informatik, 2019.
\newblock [\epfmtdoi{10.4230/LIPIcs.ITCS.2019.20}, \epfmt{arxiv}{1801.02816}]

\bibitem{CDST15}\bibhead{CDST15}
{\sc Xi~Chen, Anindya De, Rocco~A. Servedio, and Li{-}Yang Tan}: Boolean
  function monotonicity testing requires (almost) $n^{1/2}$ non-adaptive
  queries.
\newblock In {\em Proc. 47th STOC}, pp. 519--528. ACM Press, 2015.
\newblock [\epfmtdoi{10.1145/2746539.2746570}, \epfmt{arxiv}{1412.5657}]

\bibitem{CST14}\bibhead{CST14}
{\sc Xi~Chen, Rocco~A. Servedio, and Li{-}Yang Tan}: New algorithms and lower
  bounds for monotonicity testing.
\newblock In {\em Proc. 55th FOCS}, pp. 286--295. IEEE Comp. Soc. Press, 2014.
\newblock [\epfmtdoi{10.1109/FOCS.2014.38}, \epfmt{arxiv}{1412.5655}]

\bibitem{CW19}\bibhead{CW19}
{\sc Xi~Chen and Erik Waingarten}: Testing unateness nearly optimally.
\newblock In {\em Proc. 51st STOC}, pp. 547--558. ACM Press, 2019.
\newblock [\epfmtdoi{10.1145/3313276.3316351}, \epfmt{arxiv}{1904.05309}]

\bibitem{CWX17}\bibhead{CWX17}
{\sc Xi~Chen, Erik Waingarten, and Jinyu Xie}: Beyond {Talagrand} functions:
  new lower bounds for testing monotonicity and unateness.
\newblock In {\em Proc. 49th STOC}, pp. 523--536. ACM Press, 2017.
\newblock [\epfmtdoi{10.1145/3055399.3055461}, \epfmt{arxiv}{1702.06997}]

\bibitem{CWX17focs}\bibhead{CWX17focs}
{\sc Xi~Chen, Erik Waingarten, and Jinyu Xie}: Boolean unateness testing with
  $\widetilde{O}(n^{3/4})$ adaptive queries.
\newblock In {\em Proc. 58th FOCS}, pp. 868--879. IEEE Comp. Soc. Press, 2017.
\newblock [\epfmtdoi{10.1109/FOCS.2017.85}, \epfmt{arxiv}{1708.05786}]

\bibitem{DJRT13}\bibhead{DJRT13}
{\sc Kashyap Dixit, Madhav Jha, Sofya Raskhodnikova, and Abhradeep Thakurta}:
  Testing the {Lipschitz} property over product distributions with applications
  to data privacy.
\newblock In {\em Proc. 10th Theory of Cryptography Conf. (TCC'13)}, pp.
  418--436. Springer, 2013.
\newblock [\epfmtdoi{10.1007/978-3-642-36594-2\_24}, \epfmt{arxiv}{1209.4056}]

\bibitem{DRTV16}\bibhead{DRTV16}
{\sc Kashyap Dixit, Sofya Raskhodnikova, Abhradeep Thakurta, and Nithin~M.
  Varma}: Erasure-resilient property testing.
\newblock {\em SIAM J. Comput.}, 47(2):295--329, 2018.
\newblock \prelim
  \href{https://doi.org/10.4230/LIPIcs.ICALP.2016.91}{ICALP'16}.
\newblock [\epfmtdoi{10.1137/16M1075661}, \epfmt{arxiv}{1607.05786}]

\bibitem{DGLRRS99}\bibhead{DGLRRS99}
{\sc Yevgeniy Dodis, Oded Goldreich, Eric Lehman, Sofya Raskhodnikova, Dana
  Ron, and Alex Samorodnitsky}: Improved testing algorithms for monotonicity.
\newblock In {\em Proc. 3rd Internat. Workshop on Randomization and Computation
  (RANDOM'99)}, pp. 97--108. Springer, 1999.
\newblock [\epfmtdoi{10.1007/978-3-540-48413-4\_10}]

\bibitem{EKKRV00}\bibhead{EKKRV00}
{\sc Funda Erg{\"{u}}n, Sampath Kannan, Ravi Kumar, Ronitt Rubinfeld, and
  Mahesh Viswanathan}: Spot-checkers.
\newblock {\em J. Comput. System Sci.}, 60(3):717--751, 2000.
\newblock \prelim \href{https://doi.org/10.1145/276698.276757}{STOC'98}.
\newblock [\epfmtdoi{10.1006/jcss.1999.1692}]

\bibitem{Fis04}\bibhead{Fis04}
{\sc Eldar Fischer}: On the strength of comparisons in property testing.
\newblock {\em Inform. and Comput.}, 189(1):107--116, 2004.
\newblock [\epfmtdoi{10.1016/j.ic.2003.09.003}]

\bibitem{FLNRRS02}\bibhead{FLNRRS02}
{\sc Eldar Fischer, Eric Lehman, Ilan Newman, Sofya Raskhodnikova, Ronitt
  Rubinfeld, and Alex Samorodnitsky}: Monotonicity testing over general poset
  domains.
\newblock In {\em Proc. 34th STOC}, pp. 474--483. ACM Press, 2002.
\newblock [\epfmtdoi{10.1145/509907.509977}]

\bibitem{Go-book}\bibhead{Go-book}
{\sc Oded Goldreich}: {\em Introduction to Property Testing}.
\newblock Cambridge Univ. Press, 2017.
\newblock [\epfmtdoi{10.1017/9781108135252}]

\bibitem{GGLRS00}\bibhead{GGLRS00}
{\sc Oded Goldreich, Shafi Goldwasser, Eric Lehman, Dana Ron, and Alex
  Samorodnitsky}: Testing monotonicity.
\newblock {\em Combinatorica}, 20(3):301--337, 2000.
\newblock \prelim \href{http://doi.org/10.1109/SFCS.1998.743493}{FOCS'98}.
\newblock [\epfmtdoi{10.1007/s004930070011}]

\bibitem{GGR98}\bibhead{GGR98}
{\sc Oded Goldreich, Shafi Goldwasser, and Dana Ron}: Property testing and its
  connection to learning and approximation.
\newblock {\em J. ACM}, 45(4):653--750, 1998.
\newblock \prelim \href{https://doi.org/10.1109/SFCS.1996.548493}{FOCS'96}.
\newblock [\epfmtdoi{10.1145/285055.285060}]

\bibitem{GR11}\bibhead{GR11}
{\sc Oded Goldreich and Dana Ron}: On proximity-oblivious testing.
\newblock {\em SIAM J. Comput.}, 40(2):534--566, 2011.
\newblock \prelim \href{https://doi.org/10.1145/1536414.1536436}{STOC'09}.
\newblock [\epfmtdoi{10.1137/100789646}]

\bibitem{HK07}\bibhead{HK07}
{\sc Shirley Halevy and Eyal Kushilevitz}: Distribution-free property-testing.
\newblock {\em SIAM J. Comput.}, 37(4):1107--1138, 2007.
\newblock \prelim
  \href{https://doi.org/10.1007/978-3-540-45198-3\_26}{RANDOM'03}.
\newblock [\epfmtdoi{10.1137/050645804}]

\bibitem{HK08}\bibhead{HK08}
{\sc Shirley Halevy and Eyal Kushilevitz}: Testing monotonicity over graph
  products.
\newblock {\em Random Structures Algorithms}, 33(1):44--67, 2008.
\newblock \prelim
  \href{https://doi.org/10.1007/978-3-540-27836-8\_61}{ICALP'04}.
\newblock [\epfmtdoi{10.1002/rsa.20211}]

\bibitem{JR13}\bibhead{JR13}
{\sc Madhav Jha and Sofya Raskhodnikova}: Testing and reconstruction of
  {Lipschitz} functions with applications to data privacy.
\newblock {\em SIAM J. Comput.}, 42(2):700--731, 2013.
\newblock \prelim \href{https://doi.org/10.1109/FOCS.2011.13}{FOCS'11}.
\newblock [\epfmtdoi{10.1137/110840741}]

\bibitem{KLX10}\bibhead{KLX10}
{\sc Tali Kaufman, Simon Litsyn, and Ning Xie}: Breaking the
  $\epsilon$-soundness bound of the linearity test over {GF(2)}.
\newblock {\em SIAM J. Comput.}, 39(5):1988--2003, 2010.
\newblock \prelim
  \href{https://doi.org/10.1007/978-3-540-85363-3\_39}{RANDOM'08}.
\newblock [\epfmtdoi{10.1137/080715548}]

\bibitem{KMS18}\bibhead{KMS18}
{\sc Subhash Khot, Dor Minzer, and Muli Safra}: On monotonicity testing and
  {Boolean} isoperimetric-type theorems.
\newblock {\em SIAM J.\ Comput.}, 47(6):2238--2276, 2018.
\newblock \prelim \href{https://doi.org/10.1109/FOCS.2015.13}{FOCS'15}.
\newblock [\epfmtdoi{10.1137/16M1065872}]

\bibitem{KS16}\bibhead{KS16}
{\sc Subhash Khot and Igor Shinkar}: An $\widetilde{O}(n)$ queries adaptive
  tester for unateness.
\newblock In {\em Proc. 20th Internat. Workshop on Randomization and
  Computation (RANDOM'16)}, pp. 37:1--37:7. Schloss Dagstuhl--Leibniz-Zentrum
  fuer Informatik, 2016.
\newblock [\epfmtdoi{10.4230/LIPIcs.APPROX-RANDOM.2016.37},
  \epfmt{arxiv}{1608.02451}]

\bibitem{LR01}\bibhead{LR01}
{\sc Eric Lehman and Dana Ron}: On disjoint chains of subsets.
\newblock {\em J. Combin. Theory Ser. A}, 94(2):399--404, 2001.
\newblock [\epfmtdoi{10.1006/jcta.2000.3148}]

\bibitem{LW19}\bibhead{LW19}
{\sc Amit Levi and Erik Waingarten}: Lower bounds for tolerant junta and
  unateness testing via rejection sampling of graphs.
\newblock In {\em Proc. 10th Innovations in Theoret. Computer Science
  (ITCS'19)}, pp. 52:1--52:20. Schloss Dagstuhl--Leibniz-Zentrum fuer
  Informatik, 2019.
\newblock [\epfmtdoi{10.4230/LIPIcs.ITCS.2019.52}, \epfmt{arxiv}{1805.01074}]

\bibitem{PRV17}\bibhead{PRV17}
{\sc Ramesh Krishnan~S. Pallavoor, Sofya Raskhodnikova, and Nithin~M. Varma}:
  Parameterized property testing of functions.
\newblock {\em ACM Trans. Comput. Theory}, 9(4):17:1--17:19, 2018.
\newblock \prelim \href{https://doi.org/10.4230/LIPIcs.ITCS.2017.12}{ITCS'17}.
\newblock [\epfmtdoi{10.1145/3155296}]

\bibitem{PRW20}\bibhead{PRW20}
{\sc Ramesh Krishnan~S. Pallavoor, Sofya Raskhodnikova, and Erik Waingarten}:
  Approximating the distance to monotonicity of {Boolean} functions.
\newblock In {\em Proc. 31st ACM-SIAM Symp. on Discrete Algorithms (SODA'20)},
  pp. 1995--2009. ACM Press, 2020.
\newblock [\epfmtdoi{10.1137/1.9781611975994.123}, \epfmt{arxiv}{1911.06924}]

\bibitem{PRR03}\bibhead{PRR03}
{\sc Michal Parnas, Dana Ron, and Ronitt Rubinfeld}: On testing convexity and
  submodularity.
\newblock {\em SIAM J. Comput.}, 32(5):1158--1184, 2003.
\newblock \prelim \href{https://doi.org/10.1007/3-540-45726-7\_2}{RANDOM'02}.
\newblock [\epfmtdoi{10.1137/S0097539702414026}]

\bibitem{Ras99}\bibhead{Ras99}
{\sc Sofya Raskhodnikova}: Monotonicity testing.
\newblock Master's thesis, Massachusetts Institute of Technology, Cambridge,
  MA, {USA}, 1999.
\newblock \href{https://dspace.mit.edu/handle/1721.1/80562}{LINK}.

\bibitem{Ras03}\bibhead{Ras03}
{\sc Sofya Raskhodnikova}: {\em Property Testing: Theory and Applications}.
\newblock Ph.\,D.\ thesis, Massachusetts Institute of Technology, Cambridge,
  MA, {USA}, 2003.
\newblock \href{http://hdl.handle.net/1721.1/28274}{LINK}.

\bibitem{Ras10}\bibhead{Ras10}
{\sc Sofya Raskhodnikova}: Transitive-closure spanners: {A} survey.
\newblock In {\em Property Testing}, volume 6390 of {\em LNCS}, pp. 167--196.
  Springer, 2010.
\newblock [\epfmtdoi{10.1007/978-3-642-16367-8\_10}]

\bibitem{Ras16}\bibhead{Ras16}
{\sc Sofya Raskhodnikova}: Testing if an array is sorted.
\newblock In {\em Encyclopedia of Algorithms}, pp. 2219--2222. Springer, 2016.
\newblock [\epfmtdoi{10.1007/978-1-4939-2864-4\_700}]

\bibitem{RR16}\bibhead{RR16}
{\sc Sofya Raskhodnikova and Ronitt Rubinfeld}: Linearity testing/testing
  {Hadamard} codes.
\newblock In {\em Encyclopedia of Algorithms}, pp. 1107--1110. Springer, 2016.
\newblock [\epfmtdoi{10.1007/978-1-4939-2864-4\_202}]

\bibitem{RS06}\bibhead{RS06}
{\sc Sofya Raskhodnikova and Adam~D. Smith}: A note on adaptivity in testing
  properties of bounded degree graphs.
\newblock {\em Electron. Colloq. on Comput. Complexity (ECCC)}, 2006.
\newblock [\epfmt{eccc}{TR06-089}]

\bibitem{RY13}\bibhead{RY13}
{\sc Sofya Raskhodnikova and Grigory Yaroslavtsev}: Learning pseudo-{Boolean}
  $k$-{DNF} and submodular functions.
\newblock In {\em Proc. 24th ACM-SIAM Symp. on Discrete Algorithms (SODA'13)},
  pp. 1356--1368. ACM Press, 2013.
\newblock [\epfmtdoi{10.1137/1.9781611973105.98}, \epfmt{arxiv}{1208.2294}]

\bibitem{RS96}\bibhead{RS96}
{\sc Ronitt Rubinfeld and Madhu Sudan}: Robust characterizations of polynomials
  with applications to program testing.
\newblock {\em SIAM J. Comput.}, 25(2):252--271, 1996.
\newblock [\epfmtdoi{10.1137/S0097539793255151}]

\bibitem{SV14}\bibhead{SV14}
{\sc C.~Seshadhri and Jan Vondr{\'{a}}k}: Is submodularity testable?
\newblock {\em Algorithmica}, 69(1):1--25, 2014.
\newblock \prelim
  \href{https://conference.iiis.tsinghua.edu.cn/ICS2011/content/papers/21.html}{ICS'10}.
\newblock [\epfmtdoi{10.1007/s00453-012-9719-2}, \epfmt{arxiv}{1008.0831}]

\bibitem{Yao77}\bibhead{Yao77}
{\sc Andrew~Chi{-}Chih Yao}: Probabilistic computations: Toward a unified
  measure of complexity (extended abstract).
\newblock In {\em Proc. 18th FOCS}, pp. 222--227. IEEE Comp. Soc. Press, 1977.
\newblock [\epfmtdoi{10.1109/SFCS.1977.24}]

\end{thebibliography}

\end{document}
